% !TEX root = ../main.tex
\chapter{PHÂN TÍCH}

\section{Quy trình nghiệp vụ}

Quy trình nghiệp vụ của nền tảng IELTS Master English AI được xây dựng xoay quanh hai nhóm tác nhân chính: Người học (Learner) và Người quản lý (Admin). Thay vì ép người học đi theo một hành trình tuyến tính cố định, nền tảng tổ chức các chức năng thành sáu phân hệ hoàn toàn độc lập để người học tự do chọn lựa phần phù hợp với nhu cầu: phân hệ IELTS dành riêng cho người có mục tiêu luyện thi cụ thể, phân hệ Shadowing \& Dictation để rèn nghe--nói tiếng Anh tổng quát qua video YouTube, phân hệ Vocab Lab để quản lý và ôn luyện từ vựng bằng flashcard thông minh, phân hệ Community để kết nối với cộng đồng người học, phân hệ Pricing để nâng cấp quyền truy cập theo nhu cầu, và hệ thống thông báo để cập nhật các sự kiện quan trọng trong quá trình học. Cụ thể, quy trình hoạt động từng phân hệ có thể hình dung như sau:

% =================================================================
% 1. NGHIỆP VỤ NGƯỜI HỌC
% =================================================================
\subsection{Nghiệp vụ dành cho Người học}

\subsubsection{Quản lý tài khoản và Hồ sơ cá nhân}
Trước khi sử dụng bất kỳ tính năng nào, người dùng cần tạo tài khoản trên hệ thống. Quá trình đăng ký hỗ trợ hai hướng: người dùng có thể tự tạo tài khoản bằng cách cung cấp địa chỉ email và mật khẩu --- hệ thống băm mật khẩu bằng bcrypt và cấp JWT ngay sau khi đăng ký thành công --- hoặc đăng ký nhanh bằng tài khoản Google thông qua Google OAuth, trong đó hệ thống nhận ID token từ Google, xác minh theo danh sách Client ID được cấu hình (web, iOS, Android) rồi tự động tạo hồ sơ người dùng. Sau khi có tài khoản, người dùng đăng nhập và được cấp JWT để duy trì phiên làm việc trong các yêu cầu tiếp theo.

Sau khi đăng nhập thành công, người dùng có thể truy cập trang hồ sơ cá nhân --- được tổ chức thành bốn tab chức năng rõ ràng. Tab \textbf{Account Details} cho phép xem và cập nhật thông tin cá nhân như biệt danh, số điện thoại, địa chỉ và ảnh đại diện; đồng thời hiển thị thông tin gói đăng ký đang dùng bao gồm tên gói, ngày hết hạn và các hạn mức sử dụng còn lại trong tháng. Tab \textbf{Security} cho phép người dùng đổi mật khẩu bằng cách nhập mật khẩu hiện tại và mật khẩu mới; chức năng này bị vô hiệu hóa hoàn toàn đối với tài khoản đăng nhập qua Google OAuth, vì những tài khoản này không có mật khẩu trong cơ sở dữ liệu --- việc đặt mật khẩu tùy tiện có thể gây xung đột với luồng xác thực Google về sau. Tab \textbf{Gamification} là nơi theo dõi quá trình tích lũy thành tích: hệ thống cộng điểm kinh nghiệm (XP) mỗi khi người học hoàn thành bài học, đăng bài cộng đồng hay duy trì chuỗi học liên tiếp; khi đạt đủ XP ở một ngưỡng nhất định người học lên cấp và nhận huy hiệu thành tích (Achievement) tương ứng, tất cả được hiển thị tại tab này. Tab \textbf{Danger Zone} cung cấp chức năng xóa tài khoản vĩnh viễn, yêu cầu người dùng xác nhận một lần cuối trước khi thực hiện để tránh thao tác nhầm.

\subsubsection{Phân hệ IELTS}

Phân hệ IELTS là trọng tâm của toàn bộ nền tảng, được thiết kế dành riêng cho người học có nhu cầu chuẩn bị cho kỳ thi IELTS một cách có hệ thống. Khi truy cập IELTS từ thanh điều hướng chính, người học vào một không gian học tập có sidebar riêng với hai nhóm mục: nhóm \textbf{LEARNING STAGES} gồm bốn cấp độ tăng dần (Foundation, IELTS Basic, IELTS Advanced, IELTS Intensive) và nhóm \textbf{OTHER FEATURES} gồm bốn công cụ hỗ trợ (Roadmap, Statistics, Calculator, Student/Teacher). Ý tưởng thiết kế là người học bắt đầu từ Foundation để xây dựng nền tảng ngôn ngữ, tiến lên IELTS Basic để làm quen định dạng đề thi, rồi IELTS Advanced để luyện sâu theo Part, và cuối cùng là IELTS Intensive để thi thử toàn phần; trong suốt hành trình đó, các công cụ hỗ trợ giúp người học định hướng, theo dõi tiến độ và tối ưu kế hoạch ôn tập.

\paragraph{1. Foundation --- Xây dựng nền tảng ngôn ngữ}

Foundation là bước khởi đầu trong lộ trình IELTS, tập trung vào ba nền tảng cốt lõi của tiếng Anh học thuật: Từ vựng, Ngữ pháp và Phát âm.

Ở phần \textbf{Từ vựng (Vocabulary)}, nền tảng cung cấp toàn bộ loạt sách \textit{``4000 Essential English Words''} gồm 6 cuốn (Book~1 đến Book~6), mỗi cuốn gồm 600 từ vựng được phân bổ đều vào 30 unit --- tức mỗi unit có 20 từ. Người học học từng unit theo thứ tự: đầu tiên xem danh sách từ kèm định nghĩa, phiên âm và câu ví dụ, sau đó luyện qua hai bài tập trắc nghiệm (mỗi bài 20 câu), và cuối mỗi unit là một bài đọc hiểu ngắn chứa các từ vừa học kèm 5 câu hỏi kiểm tra khả năng áp dụng trong ngữ cảnh. Trong khi đọc, nếu người học bôi đen bất kỳ từ nào, hệ thống hiển thị một floating panel cho phép gửi từ đó sang Vocab Lab để Lexon AI tự động điền nội dung flashcard, giúp việc học từ vựng trong bài và ôn luyện bằng flashcard trở nên liền mạch. Người dùng FREE chỉ truy cập được hai cuốn đầu (Book~1 và Book~2); muốn học các cuốn cao hơn cần nâng cấp lên gói Premium hoặc Pro.

Ở phần \textbf{Ngữ pháp (Grammar)}, nền tảng sử dụng ba cuốn thuộc bộ \textit{Grammar in Use} --- bộ giáo trình ngữ pháp chuẩn được sử dụng rộng rãi nhất trong dạy tiếng Anh học thuật. Cuốn thứ nhất là \textit{Essential Grammar in Use} (trình độ Elementary, 115 unit, tác giả Raymond Murphy), dành cho người học đang xây dựng nền tảng ngữ pháp cơ bản. Cuốn thứ hai là \textit{English Grammar in Use} (trình độ Intermediate, 145 unit, tác giả Raymond Murphy), hướng đến người học đã có nền tảng và muốn đi sâu hơn vào các cấu trúc phức tạp. Cuốn thứ ba là \textit{Advanced Grammar in Use} (trình độ Advanced, 105 unit, tác giả Martin Hewings), dành cho người học nâng cao nhắm đến band 7.0 trở lên. Mỗi unit trong cả ba cuốn đều có cùng cấu trúc: phần trình bày lý thuyết với ví dụ minh họa cụ thể rồi đến phần bài tập thực hành (trắc nghiệm, điền vào chỗ trống, sắp xếp câu); hệ thống chấm điểm tự động và lưu tiến độ từng unit. Người dùng FREE chỉ truy cập cuốn Elementary; hai cuốn còn lại yêu cầu Premium hoặc Pro.

Ở phần \textbf{Phát âm (Pronunciation)}, hệ thống bao phủ đầy đủ \textbf{44 âm vị IPA} của tiếng Anh gồm nguyên âm đơn (monophthongs), nguyên âm đôi (diphthongs) và toàn bộ hệ thống phụ âm. Với mỗi âm, người học xem hình ảnh mô tả khẩu hình miệng, nghe audio phát âm chuẩn từ người bản ngữ, đọc mô tả kỹ thuật về vị trí lưỡi và cách tạo âm, đồng thời xem các từ ví dụ chứa âm đó ở nhiều vị trí khác nhau (đầu, giữa, cuối từ). Để luyện tập, người học nhấn nút ghi âm và đọc to từ mẫu; hệ thống gửi file âm thanh sang backend-ai, tại đó faster-whisper chuyển giọng nói thành văn bản, sau đó chuyển sang ký hiệu IPA rồi so sánh với âm chuẩn bằng thuật toán Levenshtein có trọng số --- các cặp âm dễ nhầm như /æ/--/e/ hay /l/--/r/ được gán hệ số phạt cao hơn --- kết hợp với độ tin cậy nhận dạng từ Whisper để cho ra điểm chính xác tổng hợp, giúp người học biết âm nào còn yếu để tập trung luyện lại. Người dùng FREE bị giới hạn 5 lượt ghi âm mỗi ngày; Premium và Pro không giới hạn.

\paragraph{2. IELTS Basic --- Làm quen định dạng đề thi}

Sau khi đã có nền tảng từ Foundation, người học chuyển sang IELTS Basic để làm quen với định dạng và các dạng câu hỏi chuẩn trong đề thi IELTS thực tế. IELTS Basic bao gồm cả bốn kỹ năng Listening, Reading, Writing và Speaking; mỗi kỹ năng được chia thành hai phần là \textbf{Lessons} và \textbf{Exercises}. Phần Lessons cung cấp các bài học lý thuyết giải thích từng loại câu hỏi IELTS --- ví dụ với Listening, người học được hướng dẫn cách nghe điền vào bảng biểu, cách xử lý phần nghe nhiều người nói hay cách làm Multiple Choice khi các lựa chọn rất giống nhau --- và mỗi bài học kết thúc bằng quiz củng cố ngắn để kiểm tra mức độ hiểu. Phần Exercises cho người học áp dụng lý thuyết vào các bài thực hành chuẩn IELTS; bài tập được chấm tự động và kết quả được lưu lại để theo dõi tỷ lệ chính xác theo từng dạng câu hỏi. Người dùng FREE chỉ học được tối đa 3 bài Lesson và 2 bài Exercise cho mỗi kỹ năng; khi chạm giới hạn, hệ thống hiển thị thông báo yêu cầu nâng cấp gói. Người dùng Premium và Pro không bị giới hạn.

\paragraph{3. IELTS Advanced --- Luyện sâu theo Part}

IELTS Advanced là cấp độ dành cho người học đã nắm vững lý thuyết và muốn luyện chuyên sâu theo từng Part mô phỏng sát đề thi IELTS thực tế; cấp độ này chỉ mở cho người dùng Premium trở lên, người dùng FREE sẽ thấy khóa tại đây. Với kỹ năng Listening, người học luyện qua 4 Part theo cấu trúc đề chuẩn: Part~1 (hội thoại hai người trong ngữ cảnh đời thường), Part~2 (độc thoại phi học thuật), Part~3 (thảo luận nhóm trong ngữ cảnh học thuật) và Part~4 (bài giảng hay thuyết trình học thuật). Với Reading, hệ thống cung cấp 3 Part (Part~4 không tồn tại trong IELTS Academic nên bị ẩn), mỗi Part là một đoạn văn học thuật dài kèm đầy đủ các dạng câu hỏi như True/False/Not~Given, Matching Headings, Summary Completion và Multiple Choice. Với Writing, người học luyện theo từng prompt: Task~1 mô tả biểu đồ, bảng số liệu, bản đồ hoặc sơ đồ quy trình; Task~2 viết bài luận phân tích, thảo luận hoặc phản biện. Với Speaking, người học luyện theo Part~1 (câu hỏi chủ đề quen thuộc), Part~2 (thuyết trình 2 phút từ cue card) và Part~3 (thảo luận chuyên sâu), kèm theo tab Community để xem bài tốt từ người học khác và tab My Answers để xem lại bài mình đã nộp.

Sau khi nộp bài, Listening và Reading được chấm điểm tự động ngay lập tức. Riêng Writing và Speaking cần đánh giá sâu hơn: hệ thống đưa bài vào hàng đợi RabbitMQ (\texttt{exam-grading-queue}), backend-ai nhận và gửi tới Gemini API để phân tích theo bốn tiêu chí chấm điểm IELTS chính thức --- với Writing là Task Achievement/Response, Coherence and Cohesion, Lexical Resource và Grammatical Range and Accuracy; với Speaking là Fluency and Coherence, Lexical Resource, Grammatical Range and Accuracy và Pronunciation --- rồi trả về band điểm ước tính (thang 0--9) cho từng tiêu chí kèm nhận xét chi tiết và gợi ý cải thiện cụ thể. Sau khi chấm xong, backend-ai gọi HTTP callback về backend-core để cập nhật kết quả và gửi thông báo cho người học. Người dùng Premium được sử dụng tối đa 10 lượt chấm AI Writing và 10 lượt chấm AI Speaking mỗi tháng; người dùng Pro không giới hạn.

\paragraph{4. IELTS Intensive --- Thi thử toàn phần}

IELTS Intensive cung cấp trải nghiệm thi thử sát thực nhất với hai mục: \textbf{Mock Tests} và \textbf{Test History}. Khi vào Mock Tests, người học chọn đề thi thử và bắt đầu phiên thi với đồng hồ đếm ngược theo đúng thời gian quy định của IELTS thực tế (Listening 40 phút, Reading 60 phút, Writing 60 phút, Speaking mô phỏng 3 Part với AI). Trong suốt phiên thi, nếu người học muốn thoát, hệ thống cảnh báo và lưu lại trạng thái phiên. Sau khi nộp toàn bộ, phiên thi chuyển trạng thái lần lượt: \texttt{SUBMITTED} khi người học nộp xong, \texttt{GRADING} trong khi backend-ai xử lý phần Writing và Speaking, và \texttt{GRADED} khi kết quả đã đầy đủ; nếu xảy ra sự cố trong quá trình chấm, phiên chuyển sang \texttt{GRADING\_FAILED} và người học được thông báo. Kết quả cuối cùng tổng hợp band điểm theo từng kỹ năng và band tổng thể theo quy tắc làm tròn chính thức của IELTS. Tại Test History, người học xem lại danh sách toàn bộ phiên thi đã thực hiện, lọc theo kỹ năng hoặc thời gian, nhấn vào từng phiên để xem chi tiết đáp án và nhận xét, đồng thời có thể xóa phiên không cần thiết. Người dùng FREE chỉ lưu tối đa 3 phiên thi gần nhất; khi có phiên thứ 4, phiên cũ nhất tự động bị ghi đè. Người dùng Premium và Pro lưu không giới hạn.

\paragraph{5. Roadmap --- Lộ trình học cá nhân hóa}

Roadmap giúp người học có một kế hoạch học tập cụ thể thay vì tự mày mò chọn bài. Để sử dụng, người học hoàn thành bước onboarding gồm các câu hỏi: band mục tiêu (ví dụ 6.5 hay 7.0), mức độ cam kết học mỗi ngày (30 phút, 1 giờ hay 2 giờ), và ngày thi dự kiến. Hệ thống gợi ý người học làm thêm bài kiểm tra chẩn đoán (Diagnostic Quiz) cho ba kỹ năng Listening, Reading và Writing để đánh giá điểm xuất phát; người học có thể bỏ qua nếu đã biết rõ trình độ của mình. Từ tất cả thông tin đó, hệ thống sinh ra lộ trình học tập theo tuần: mỗi tuần có danh sách bài học và bài tập cần hoàn thành từ Foundation, IELTS Basic và IELTS Advanced phù hợp với trình độ hiện tại. Mỗi bước trong lộ trình tự động được đánh dấu ``Hoàn thành'' khi người học đạt từ 90\% trở lên trong bài kiểm tra tương ứng.

\paragraph{6. Statistics --- Theo dõi tiến độ tổng thể}

Statistics cung cấp cái nhìn toàn diện về hành trình học tập qua năm tab thống kê chuyên biệt. Tab \textbf{Overview} là trang tổng hợp: hiển thị band điểm ước tính hiện tại cho từng kỹ năng (tính từ lịch sử bài làm gần nhất), biểu đồ heatmap hoạt động học tập theo ngày trong năm (tương tự GitHub contribution graph), tiến độ hoàn thành mục tiêu học hàng ngày, và bộ đếm ngược số ngày còn đến ngày thi đã nhập trong Roadmap. Tab \textbf{Foundation} chi tiết tiến độ học Vocabulary (số unit đã hoàn thành trong từng cuốn), Grammar (số unit đã học và điểm trung bình theo từng cuốn) và Pronunciation (tỷ lệ chính xác theo từng nhóm âm). Tab \textbf{Basic} thống kê số bài Lesson và Exercise đã hoàn thành cùng điểm trung bình theo từng kỹ năng, giúp người học nhận ra kỹ năng nào còn đang yếu. Tab \textbf{Advanced} tổng kết các phiên luyện theo Part và kỹ năng cùng xu hướng cải thiện theo thời gian. Tab \textbf{Intensive} thể hiện lịch sử band điểm các phiên thi thử dưới dạng đường biểu diễn, giúp người học thấy rõ mình đang tiến bộ như thế nào qua từng lần thi thử.

\paragraph{7. Calculator --- Quy đổi điểm IELTS}

Calculator là công cụ tiện ích giúp người học hiểu rõ hệ thống tính điểm IELTS và tự dự đoán kết quả. Công cụ tổ chức theo bốn tab kỹ năng: tab Listening và tab Reading cho phép nhập số câu trả lời đúng (0--40 câu) rồi tra cứu band điểm tương ứng theo bảng quy đổi chính thức của IELTS; tab Writing và tab Speaking hiển thị bảng descriptors chi tiết mô tả từng tiêu chí từ band 1.0 đến 9.0, giúp người học tự đánh giá bài viết hay bài nói của mình đang ở mức nào. Ngoài bốn tab kỹ năng, Calculator còn có mục \textbf{Overall Band}: người học nhập band của cả bốn kỹ năng, hệ thống tính band tổng theo quy tắc làm tròn chính thức (cộng bốn band chia bốn rồi làm tròn đến 0.5 gần nhất theo chiều có lợi hơn cho người thi).

\paragraph{8. Student/Teacher --- Kết nối học viên và giáo viên}

Student/Teacher là tính năng hỗ trợ mô hình dạy--học có giám sát, nơi người học và giáo viên liên kết tài khoản để giáo viên theo dõi tiến độ của học viên mình phụ trách. Giao diện gồm hai tab tùy theo vai trò. Khi đóng vai học viên, người học mở tab \textbf{Student}, nhập Teacher ID của giáo viên muốn liên kết và gửi yêu cầu kết nối; hệ thống lưu mối liên kết vào cơ sở dữ liệu và người học có thể xem lại danh sách giáo viên đang theo dõi mình, hủy liên kết bất kỳ lúc nào nếu không cần nữa. Khi đóng vai giáo viên, người học mở tab \textbf{Teacher}, hệ thống hiển thị User ID của chính mình để copy và chia sẻ cho học viên; sau khi học viên liên kết, giáo viên thấy danh sách đầy đủ các học viên kèm thông tin tiến độ học chi tiết của từng người --- band điểm các kỳ thi gần nhất, số bài hoàn thành trong tuần và lịch sử hoạt động --- từ đó đưa ra nhận xét và định hướng kịp thời.

\subsubsection{Phân hệ Shadowing và Dictation}

Shadowing \& Dictation là phân hệ hoàn toàn độc lập với IELTS, phục vụ mục đích luyện kỹ năng nghe và nói tiếng Anh tổng quát thông qua các video YouTube đa dạng chủ đề --- từ tin tức, podcast, phim tài liệu đến các cuộc trò chuyện đời thường. Phân hệ phù hợp cho bất kỳ người học tiếng Anh nào, không nhất thiết phải luyện thi IELTS, và gồm ba khu vực: thư viện Shadowing, thư viện Dictation do Admin biên soạn, và khu vực My Videos dành cho người học muốn tự nhập video riêng.

\paragraph{1. Shadowing --- Luyện phát âm bằng nói đuổi}

Shadowing là kỹ thuật luyện nói bằng cách theo dõi một đoạn audio và nói đuổi theo gần như đồng thời, giúp cải thiện phát âm, ngữ điệu, nhịp điệu và độ lưu loát. Tại thư viện Shadowing, người học xem danh sách bài do Admin xuất bản và lọc theo danh mục chủ đề hoặc theo trạng thái cá nhân (Not started, In progress, Completed) để tập trung vào bài chưa hoàn thành. Khi vào bài luyện, giao diện hiển thị video YouTube phát qua iframe bên cạnh danh sách transcript được tô sáng từng câu đồng bộ theo tiến độ video. Người học nghe câu, dừng video và nói lại câu đó; có thể ghi âm để nghe lại giọng mình và so sánh với audio gốc. Sau khi hoàn thành từng câu, hệ thống lưu tiến độ câu theo câu vào cơ sở dữ liệu, tính phần trăm hoàn thành bài, và nếu người học luyện đủ trong ngày thì cập nhật chuỗi học liên tiếp (streak). Người dùng FREE chỉ truy cập tối đa 5 bài trong thư viện; muốn học nhiều hơn cần nâng cấp gói.

\paragraph{2. Dictation --- Luyện nghe bằng điền từ}

Dictation luyện kỹ năng nghe theo phương pháp điền từ vào chỗ trống: người học nghe từng câu audio và điền lại các từ bị ẩn. Tại thư viện Dictation, người học chọn bài từ danh sách video do Admin chuẩn bị transcript sẵn. Khi vào bài, người học trước tiên chọn \textbf{cấp độ khó} phù hợp với trình độ: Beginner (70\% số từ được hiển thị, chỉ ẩn 30\%), Intermediate (50\% hiển thị), Advanced (30\% hiển thị) và Expert (toàn bộ từ bị ẩn, phải nghe và điền hoàn toàn). Trong khi nghe, nếu gặp từ khó, người học có thể dùng tính năng \textbf{gợi ý tiệm tiến}: nhấn lần đầu thì hệ thống hiện chữ cái đầu tiên của từ cần điền, nhấn lần hai thì hiện thêm chữ cái cuối, nhấn lần ba thì hiện toàn bộ từ --- cơ chế này khuyến khích người học cố gắng trước khi xem đáp án thay vì bỏ cuộc ngay. Sau khi điền xong một câu, người học nhấn kiểm tra: ô đúng chuyển sang màu xanh, ô sai chuyển sang màu đỏ; với ô sai, người học nghe lại câu và điền lại cho đến khi đúng rồi mới được chuyển sang câu tiếp theo. Kết quả hoàn thành bài được lưu lại và streak được cập nhật. Người dùng FREE được truy cập tối đa 5 bài Dictation trong thư viện.

\paragraph{3. My Shadowing và My Dictation --- Tự nhập video YouTube}

Đối với người dùng Premium và Pro, phân hệ còn cung cấp thêm khu vực My Videos, nơi người học tự xây dựng thư viện luyện tập cá nhân từ bất kỳ video YouTube công khai nào. Để thêm bài mới, người học dán URL video YouTube vào ô nhập liệu và xác nhận; hệ thống tạo một bản ghi với trạng thái \texttt{PROCESSING} và đưa tác vụ vào hàng đợi RabbitMQ (\texttt{dictation-transcription-queue}). Ở phía backend-ai, TranscriptionConsumer nhận tác vụ, tải audio từ YouTube về, dùng mô hình faster-whisper để chuyển đổi giọng nói thành văn bản có mốc thời gian từng từ, tự động phân tách thành các câu dựa trên dấu câu và ngữ điệu, rồi cập nhật trạng thái video sang \texttt{READY}. Trong thời gian đó, giao diện tự động polling mỗi 5 giây để kiểm tra; khi nhận được \texttt{READY}, giao diện chuyển sang hiển thị bài luyện đầy đủ và người học có thể bắt đầu sử dụng y hệt như bài trong thư viện chuẩn --- Shadowing thì nói đuổi theo transcript, Dictation thì điền từ vào chỗ trống. Người học cũng có thể xóa video không dùng nữa khỏi thư viện cá nhân. Tính năng này bị khóa hoàn toàn với người dùng FREE.

\subsubsection{Phân hệ Community}

Community là phân hệ mạng xã hội học tập nội bộ, nơi người học chia sẻ kinh nghiệm, hỏi đáp, đăng kết quả thi để động viên nhau và theo dõi thứ hạng trong cộng đồng. Phân hệ được truy cập từ thanh điều hướng chính và có sidebar riêng gồm ba mục: Feed, My Activity và Leaderboard.

\paragraph{1. Feed --- Dòng thông tin cộng đồng}

Feed là trang chủ của Community, hiển thị các bài đăng từ toàn bộ người học theo thứ tự mới nhất trước. Người học có thể lọc bài theo ba loại: \textbf{All Posts} (tất cả bài đăng), \textbf{Study Tips} (mẹo và kinh nghiệm học tập) và \textbf{Achievements} (khoe kết quả, cột mốc đạt được). Để đăng bài mới, người học nhấn vào nút tạo bài và một modal soạn thảo xuất hiện: nhập tiêu đề, viết nội dung, gán tag phân loại để bài dễ tìm kiếm, và đính kèm tối đa 4 ảnh minh họa (ảnh được tải lên Cloudinary và URL được nhúng vào bài). Ngay sau khi đăng thành công, bài xuất hiện ở đầu Feed và người đăng nhận +5 XP; nếu đây là lần đầu tiên đăng bài trên nền tảng, hệ thống cũng cấp thêm huy hiệu ``First Post''. Với mỗi bài đăng, các thành viên khác có thể tương tác qua bốn hình thức: nhấn Like để ghi nhận bài hay, nhấn Bookmark để lưu đọc sau, bình luận trực tiếp bên dưới, hoặc reply vào bình luận của người khác để tạo ra chuỗi thảo luận.

\paragraph{2. My Activity --- Hoạt động cá nhân}

My Activity là trang lưu trữ cá nhân trong phạm vi Community. Về mặt kỹ thuật, đây là Feed nhưng được lọc chỉ hiển thị các bài đăng mà \texttt{authorId} trùng với ID của người học đang đăng nhập. Tại đây, người học xem lại toàn bộ lịch sử bài đăng của mình, thấy số lượt Like và Comment mà từng bài nhận được, và có thể xóa bài nếu cần.

\paragraph{3. Leaderboard --- Bảng xếp hạng cộng đồng}

Leaderboard tạo yếu tố thi đua lành mạnh trong cộng đồng, hiển thị hai bảng xếp hạng song song với top 10 người dùng. Bảng \textbf{XP tuần này} xếp hạng dựa trên tổng điểm kinh nghiệm (\texttt{xp\_weekly}) tích lũy trong tuần hiện tại từ tất cả hoạt động học tập và tương tác cộng đồng. Bảng \textbf{Streak} xếp hạng dựa trên chuỗi ngày học liên tiếp hiện tại, tức số ngày liên tiếp gần nhất mà người học có ít nhất một hoạt động học tập. Mỗi hàng trong bảng hiển thị thứ hạng, ảnh đại diện, tên người dùng và chỉ số tương ứng; nếu người học đang đăng nhập có mặt trong bảng, dòng đó được làm nổi bật bằng đường viền bên trái màu chủ đạo để dễ nhận ra vị trí của mình.

\subsubsection{Phân hệ Vocab Lab}

Vocab Lab là phân hệ flashcard cá nhân hóa, không giới hạn trong phạm vi IELTS mà phục vụ việc học và ghi nhớ bất kỳ loại từ vựng hay kiến thức nào theo phương pháp khoa học. Nền tảng ứng dụng hai kỹ thuật ghi nhớ hiện đại: \textbf{Active Recall} (nhớ chủ động --- buộc não phải tự tái hiện thông tin thay vì chỉ đọc lại) và \textbf{Spaced Repetition} (lặp lại ngắt quãng --- ôn lại đúng lúc não sắp quên để gia cố trí nhớ dài hạn). Cốt lõi của hệ thống Spaced Repetition là thuật toán FSRS (Free Spaced Repetition Scheduler), tính khoảng thời gian ôn tập tiếp theo dựa trên độ bền trí nhớ (Stability $S$) và tỉ lệ nhớ mục tiêu theo công thức~(\ref{eq:fsrs-interval}):
\begin{equation}
    I' = S \times \dfrac{\ln(R_{target})}{\ln(0{,}9)}
    \label{eq:fsrs-interval}
\end{equation}
Trong đó $I'$ là số ngày đến lần ôn tiếp theo, $S$ là độ bền trí nhớ hiện tại của thẻ (tăng dần theo mỗi lần ôn thành công), $R_{target}$ là tỉ lệ nhớ mục tiêu (mặc định 0{,}9, tương ứng 90\%), và mẫu số $\ln(0{,}9)$ phản ánh đường cong quên tự nhiên theo lý thuyết Ebbinghaus. Nhờ cơ chế này, thẻ dễ sẽ được ôn thưa hơn và thẻ khó sẽ được ôn thường xuyên hơn, tối ưu thời gian học mà không lãng phí công sức vào những gì đã nhớ chắc.

Trong mỗi \textbf{phiên ôn tập (Study Session)}, người học lần lượt xem các thẻ --- gồm thẻ mới (số lượng do người học tự đặt mỗi ngày) và thẻ đến lịch ôn theo FSRS --- và sau khi trả lời, đánh giá mức độ ghi nhớ qua bốn nút:

\begin{itemize}
    \item \textbf{Again:} người học quên hoàn toàn. Hệ thống đặt lại thẻ về trạng thái Relearning và giảm mạnh Stability; thẻ xuất hiện lại sau vài phút ngay trong phiên ôn hiện tại.
    \item \textbf{Hard:} người học nhớ nhưng còn vất vả. FSRS tính khoảng cách ngắn hơn mức tối ưu và điều chỉnh giảm nhẹ Difficulty của thẻ.
    \item \textbf{Good:} người học nhớ đúng như kỳ vọng. FSRS áp dụng công thức~(\ref{eq:fsrs-interval}) để tính interval bình thường; Stability tăng đều.
    \item \textbf{Easy:} người học nhớ rất chắc, vượt kỳ vọng. FSRS áp dụng hệ số easy bonus, interval tăng nhanh hơn đáng kể; Difficulty giảm nhẹ giúp thẻ xuất hiện thưa hơn về lâu dài.
\end{itemize}

Vocab Lab được truy cập qua nút riêng biệt trên thanh điều hướng chính, có badge đỏ hiển thị số thẻ cần ôn hôm nay để nhắc nhở người học. Bên trong là năm tab chức năng chính.

\paragraph{1. Decks --- Quản lý bộ thẻ}

Decks là trang quản lý tổng thể các bộ flashcard. Người học có thể tạo bộ thẻ mới bằng cách đặt tên và chọn màu đại diện, đổi tên hoặc xóa bộ thẻ đã có, và nhập bộ thẻ từ bên ngoài theo định dạng file \texttt{.lexon} (một định dạng JSON tùy chỉnh lưu cấu trúc card type, field và toàn bộ nội dung thẻ). Mỗi bộ thẻ được hiển thị kèm ba con số trạng thái quan trọng: \textbf{New} (thẻ chưa bao giờ ôn, chờ được đưa vào học lần đầu), \textbf{Learning} (thẻ đang trong giai đoạn học ban đầu, interval còn rất ngắn) và \textbf{Due} (thẻ đã đến lịch ôn hôm nay theo thuật toán FSRS). Người học nhấn nút Study trên một bộ thẻ để bắt đầu phiên ôn tập. Ngoài ra, người học cũng có thể xuất bộ thẻ ra file \texttt{.lexon} để chia sẻ hoặc sao lưu, và đăng bộ thẻ lên Community Marketplace. Người dùng FREE giới hạn 3 bộ thẻ và tối đa 50 thẻ trong mỗi bộ; Premium và Pro không giới hạn.

\paragraph{2. Add --- Tạo flashcard mới}

Add là giao diện tạo flashcard mới với hệ thống Card Type linh hoạt. Mỗi Card Type do người học tự định nghĩa gồm: danh sách Fields (ví dụ: Từ vựng, Nghĩa tiếng Việt, Phát âm IPA, Câu ví dụ) với kiểu dữ liệu là văn bản hoặc đa phương tiện, và Template quy định mặt trước và mặt sau của thẻ sẽ hiển thị field nào. Hệ thống cung cấp sẵn Card Type \textit{Basic} (Field Front và Field Back dạng văn bản) cho những ai không muốn tùy chỉnh sâu. Khi tạo thẻ, người học chọn bộ thẻ đích, chọn Card Type, điền nội dung cho từng field và gán tags để tìm kiếm sau. Điểm nổi bật nhất là tính năng \textbf{Lexon AI prefill}: khi người học bôi đen một từ hay cụm từ trong bất kỳ phần nào của hệ thống --- Foundation Vocabulary, IELTS Basic, IELTS Advanced --- một floating panel hiện ra cho phép chọn ``Thêm vào Vocab Lab''; Lexon AI đọc từ đã chọn và ngữ cảnh xung quanh, tự động sinh nội dung cho tất cả các field của Card Type đang dùng rồi điền sẵn vào tab Add để người học xem lại, chỉnh sửa nếu cần và lưu. Người dùng FREE không thể dùng tính năng AI sinh thẻ; Premium được tối đa 50 lượt mỗi tháng, Pro không giới hạn.

\paragraph{3. Browse --- Duyệt và chỉnh sửa thẻ}

Browse là trình duyệt toàn bộ flashcard với khả năng lọc và chỉnh sửa nhanh. Sidebar bên trái cho phép lọc theo bộ thẻ cụ thể, theo trạng thái thẻ trong chu kỳ FSRS (New, Learning, Review, Relearning) hoặc theo tags đã gán. Danh sách thẻ hỗ trợ điều hướng bằng phím ↑/↓ để duyệt nhanh mà không cần dùng chuột. Khi chọn một thẻ, phần chi tiết bên phải cho phép chỉnh sửa nội dung từng field trực tiếp, upload ảnh hoặc file audio vào field đa phương tiện, và xem các thông số FSRS hiện tại của thẻ như Stability, Difficulty và ngày ôn tiếp theo.

\paragraph{4. Stats --- Thống kê SRS}

Stats cung cấp nhiều biểu đồ giúp người học hiểu rõ tình trạng bộ thẻ và thói quen ôn tập của mình. Biểu đồ donut hiển thị phân bố thẻ theo bốn trạng thái SRS (New, Learning, Review, Relearning). Biểu đồ thanh theo ngày thể hiện lịch sử số thẻ đã ôn trong 30 ngày qua. Biểu đồ dự báo (Forecast Chart) cho thấy số thẻ dự kiến đến lịch ôn trong 30 ngày tới, giúp người học lên kế hoạch để không bị dồn quá nhiều thẻ vào một ngày. Biểu đồ hoạt động theo giờ (Hourly Activity Chart) cho thấy người học thường ôn tập vào khung giờ nào nhiều nhất trong ngày.

\paragraph{5. Community Marketplace --- Kho thẻ cộng đồng}

Community Marketplace là nơi người học chia sẻ và tải về bộ flashcard từ cộng đồng. Tab \textbf{Explore} hiển thị bộ thẻ đã được người học khác xuất bản, phân loại theo danh mục chủ đề (English, IELTS, TOEFL, TOEIC, Academic và nhiều danh mục khác), có thể sắp xếp theo phổ biến nhất hoặc mới nhất, và tìm kiếm theo từ khóa. Khi tìm được bộ thẻ phù hợp, người học nhấn nút Nhập để sao chép toàn bộ nội dung về tài khoản của mình dưới dạng một bộ thẻ mới. Tab \textbf{My Published} hiển thị các bộ thẻ mà người học đã xuất bản kèm số lượt nhập từ cộng đồng. Người dùng FREE không thể nhập bộ thẻ từ Marketplace và cũng không thể xuất bản bộ thẻ của mình; tính năng này chỉ dành cho Premium và Pro.

\subsubsection{Phân hệ Pricing và Subscription}

Nền tảng cung cấp ba gói thành viên với quyền hạn tăng dần nhằm đáp ứng nhu cầu học tập khác nhau. Gói \textbf{FREE} miễn phí hoàn toàn nhưng giới hạn hầu hết tính năng nâng cao --- phù hợp để người dùng trải nghiệm trước khi quyết định đầu tư. Gói \textbf{PREMIUM} mở khóa đầy đủ các tính năng học cốt lõi với giá \$9{,}99 mỗi tháng hoặc \$79{,}99 mỗi năm (tiết kiệm khoảng 33\% so với thanh toán theo tháng), đây là gói phù hợp cho phần lớn người học đang luyện thi IELTS. Gói \textbf{PRO} bổ sung thêm hạn mức không giới hạn ở mọi tính năng với giá \$19{,}99 mỗi tháng hoặc \$159{,}99 mỗi năm, dành cho người học cường độ cao hoặc giáo viên cần quản lý nhiều học viên. Người dùng mới đăng ký lần đầu đều đủ điều kiện nhận \textbf{7 ngày dùng thử Premium miễn phí}, mỗi tài khoản chỉ được một lần duy nhất. Bảng~\ref{tab:feature-limits} tóm tắt các giới hạn cụ thể:

\begin{table}[H]
\centering
\caption{Giới hạn tính năng theo gói thành viên}
\label{tab:feature-limits}
\begin{tabular}{|l|c|c|c|}
\hline
\textbf{Tính năng} & \textbf{FREE} & \textbf{PREMIUM} & \textbf{PRO} \\
\hline
Foundation Vocabulary & 2 cuốn đầu & Tất cả 6 cuốn & Tất cả 6 cuốn \\
Foundation Grammar & Chỉ Elementary & Cả 3 cấp độ & Cả 3 cấp độ \\
Pronunciation (lượt/ngày) & 5 & Không giới hạn & Không giới hạn \\
IELTS Basic lesson/kỹ năng & 3 & Không giới hạn & Không giới hạn \\
IELTS Basic exercise/kỹ năng & 2 & Không giới hạn & Không giới hạn \\
IELTS Advanced & Không & Có & Có \\
AI chấm Writing (lượt/tháng) & 0 & 10 & Không giới hạn \\
AI chấm Speaking (lượt/tháng) & 0 & 10 & Không giới hạn \\
Lịch sử bài thi Intensive & 3 bài gần nhất & Không giới hạn & Không giới hạn \\
Shadowing/Dictation thư viện & 5 bài/loại & Không giới hạn & Không giới hạn \\
Nhập YouTube (My videos) & Không & Có & Có \\
Vocab Lab: số bộ thẻ & 3 & Không giới hạn & Không giới hạn \\
Vocab Lab: thẻ/bộ thẻ & 50 & Không giới hạn & Không giới hạn \\
Lexon AI sinh thẻ (lượt/tháng) & 0 & 50 & Không giới hạn \\
Marketplace nhập/xuất bản & Không & Có & Có \\
\hline
\end{tabular}
\end{table}

Để nâng cấp, người học truy cập trang Pricing, xem chi tiết từng gói và nhấn nút đăng ký. Thanh toán được thực hiện qua cổng \textbf{VNPay} với chữ ký HMAC-SHA512 đảm bảo toàn vẹn dữ liệu; sau khi hoàn tất thanh toán, người học được điều hướng về trang xác nhận và gói thành viên kích hoạt ngay lập tức. Để đảm bảo người học không bị gián đoạn đột ngột, hệ thống cron chạy hàng ngày lúc 2:00 sáng để gửi email nhắc nhở sắp hết hạn vào các mốc 7 ngày, 3 ngày và 1 ngày trước; sau khi hết hạn, tài khoản còn thêm 3 ngày gia hạn trước khi tự động downgrade về gói FREE. Phiên dùng thử trial cũng kết thúc đúng hạn theo cùng lịch cron này.

\subsubsection{Xem thông báo}

Hệ thống thông báo giúp người học không bỏ lỡ bất kỳ sự kiện quan trọng nào trong hành trình học tập. Thông báo được gửi tự động trong các trường hợp: kết quả chấm điểm Writing hoặc Speaking từ AI đã trả về (sau khi backend-ai hoàn thành và gọi HTTP callback về backend-core), hoàn thành một bài Dictation và cập nhật streak, đạt một cột mốc chuỗi học liên tiếp mới (ví dụ đạt 7 ngày, 30 ngày hay 100 ngày streak), nâng cấp hoặc gia hạn gói thành viên thành công, sắp hết hạn gói theo lịch cron 7/3/1 ngày trước, và các thông báo từ Admin gửi đến toàn bộ người dùng hay một nhóm cụ thể. Người học nhấn vào biểu tượng chuông trên thanh điều hướng để mở danh sách thông báo theo thứ tự mới nhất; có thể đánh dấu đã đọc từng thông báo hoặc tất cả cùng lúc, và nhấn vào bất kỳ thông báo nào để được điều hướng thẳng đến trang liên quan.

% =================================================================
% 2. NGHIỆP VỤ NGƯỜI QUẢN LÝ
% =================================================================
\subsection{Nghiệp vụ dành cho Người quản lý}

\subsubsection{Quản lý nội dung Shadowing và Dictation}

Người quản lý chịu trách nhiệm duy trì và mở rộng thư viện bài học cho cả hai phân hệ Shadowing và Dictation thông qua giao diện quản trị. Để tạo bài Shadowing mới, người quản lý nhập các thông tin cần thiết bao gồm tiêu đề bài học, ID video YouTube, thư mục phân loại, danh mục chủ đề, thời lượng video và quan trọng nhất là danh sách transcript đã được biên soạn kỹ lưỡng: mỗi câu transcript gắn kèm mốc thời gian bắt đầu và kết thúc (tính bằng mili giây) để hệ thống đồng bộ tô sáng transcript theo tiến độ video khi người học luyện tập. Sau khi điền đầy đủ và kiểm tra, người quản lý xuất bản bài học và bài hiển thị ngay trong thư viện của người học.

Để tạo bài Dictation, người quản lý nhập ID video YouTube tương tự, đồng thời cung cấp transcript chuẩn xác đến từng từ --- đây là ground-truth mà hệ thống dùng để so khớp đáp án khi người học điền từ vào chỗ trống; một từ sai chính tả hay thiếu dấu câu cũng có thể ảnh hưởng đến việc chấm điểm nên người quản lý cần kiểm tra kỹ trước khi xuất bản. Ngoài ra, người quản lý có thể chỉnh sửa hoặc xóa các bài học đã tạo khi cần cập nhật nội dung hoặc gỡ bỏ bài không còn phù hợp; thay đổi có hiệu lực ngay lập tức đối với tất cả người học.

\subsubsection{Quản lý gói đăng ký thành viên}

Người quản lý có quyền truy cập trang quản lý subscription để nắm tổng quan tình hình người dùng theo từng gói. Trang hiển thị số lượng tài khoản đang sử dụng gói FREE, PREMIUM và PRO cùng với tổng doanh thu theo thời gian. Danh sách chi tiết cho phép lọc theo gói thành viên để xem trạng thái đăng ký của từng học viên: ngày bắt đầu, ngày hết hạn và trạng thái hiện tại (đang hoạt động, hết hạn hay đang trong giai đoạn gia hạn). Khi có yêu cầu đặc biệt --- chẳng hạn cấp gói thử nghiệm cho đối tác, bồi thường cho người dùng gặp sự cố thanh toán, hay hỗ trợ người dùng theo chương trình khuyến mãi --- người quản lý có thể cấp gói thành viên thủ công cho một tài khoản cụ thể với thời hạn tùy chỉnh, không cần người dùng thực hiện thanh toán qua VNPay.

\section{Use-case tổng quát}
Hình~\ref{fig:usecase-overview} thể hiện sơ đồ use-case tổng quát của hệ thống, phân theo hai nhóm tác nhân chính là Người học và Người quản lý.
\begin{figure}[H]
    \centering
    \includegraphics[width=\linewidth]{use-case.png}
    \caption{Sơ đồ use-case tổng quát của hệ thống luyện thi IELTS}
    \label{fig:usecase-overview}
\end{figure}

\section{Danh sách tác nhân và mô tả}

Hệ thống có hai tác nhân chính tương tác trực tiếp, được mô tả trong bảng~\ref{tab:actors}:

\begin{table}[H]
    \centering
    \renewcommand{\arraystretch}{1.3}
    \begin{tabularx}{\linewidth}{|l|X|}
        \hline
        \textbf{Tác nhân} & \textbf{Mô tả} \\ \hline
        Người học (Learner) &
        Là đối tượng sử dụng chính của hệ thống, bao gồm bất kỳ ai có nhu cầu học tiếng Anh hoặc luyện thi IELTS. Người học đăng ký tài khoản bằng email/mật khẩu hoặc Google OAuth, sau đó truy cập sáu phân hệ hoàn toàn độc lập theo nhu cầu riêng. Trong phân hệ \textbf{IELTS}, người học xây dựng nền tảng ngôn ngữ qua từ vựng (6 cuốn \textit{4000 Essential English Words}), ngữ pháp (3 cấp độ \textit{Grammar in Use}) và phát âm (44 âm IPA); luyện lý thuyết dạng câu hỏi và bài tập chuẩn ở IELTS Basic; luyện sâu theo từng Part ở IELTS Advanced (chỉ Premium/Pro); thi thử toàn phần và xem lịch sử thi ở IELTS Intensive; đồng thời sử dụng Roadmap, Statistics, Calculator và Student/Teacher để định hướng và theo dõi tiến độ. Trong phân hệ \textbf{Shadowing \& Dictation}, người học luyện nói đuổi (Shadowing) và nghe chép chính tả (Dictation) từ thư viện video YouTube do Admin biên soạn; người dùng Premium/Pro có thể tự nhập video YouTube riêng để luyện. Trong phân hệ \textbf{Community}, người học đăng bài chia sẻ kinh nghiệm, tương tác với bài của người khác qua Like/Comment, theo dõi hoạt động cá nhân và xem bảng xếp hạng XP tuần và Streak. Trong phân hệ \textbf{Vocab Lab}, người học tạo và ôn luyện bộ flashcard cá nhân với thuật toán FSRS, dùng Lexon AI sinh thẻ từ nội dung bài học, duyệt và chỉnh sửa thẻ qua Browse, xem thống kê SRS, và trao đổi bộ thẻ qua Community Marketplace. Trong phân hệ \textbf{Pricing}, người học đăng ký gói thành viên (FREE/PREMIUM/PRO) và thanh toán qua VNPay. Ngoài ra, người học nhận và quản lý \textbf{thông báo} về kết quả chấm điểm, streak và trạng thái gói. Người học cũng có thể mở rộng vai trò thành giáo viên để theo dõi tiến độ học viên đã liên kết qua tính năng Student/Teacher. \\ \hline
        Người quản lý (Admin) &
        Là người có vai trò quản lý nội dung và vận hành hệ thống, truy cập qua giao diện Admin riêng biệt. Người quản lý tạo, chỉnh sửa và xóa bài học Shadowing và Dictation cho thư viện chuẩn (nhập ID video YouTube, transcript có timestamp cho Shadowing; transcript ground-truth cho Dictation); quản lý danh sách đăng ký thành viên (theo dõi trạng thái từng tài khoản theo gói FREE/PREMIUM/PRO, xem doanh thu tổng hợp); và cấp gói thành viên thủ công cho học viên cụ thể khi cần. \\ \hline
    \end{tabularx}
    \caption{Danh sách tác nhân và mô tả}
    \label{tab:actors}
\end{table}


\section{Danh sách các tình huống hoạt động chính}
Bảng~\ref{tab:uc_list} liệt kê toàn bộ các use case của hệ thống kèm mô tả ngắn gọn về chức năng và tác nhân liên quan.

\renewcommand{\arraystretch}{1.3}
% 3 cột p{}, divisor=150: (1.5+4.5+9.0)×10
\setlength{\tbu}{\dimexpr(\linewidth - 6\tabcolsep - 4\arrayrulewidth)/150\relax}
\begin{longtable}{|p{15\tbu}|p{45\tbu}|p{90\tbu}|}
    \hline
    \textbf{ID} & \textbf{Tên use-case} & \textbf{Mô tả} \\ \hline
    \endfirsthead
    \multicolumn{3}{r}{\textit{(tiếp trang trước)}} \\
    \hline
    \textbf{ID} & \textbf{Tên use-case} & \textbf{Mô tả} \\ \hline
    \endhead
    \hline \multicolumn{3}{r}{\textit{(tiếp trang sau)}} \\
    \endfoot
    \endlastfoot
    UC01 & Đăng ký tài khoản & Người học tạo tài khoản mới bằng email và mật khẩu (hệ thống băm mật khẩu bằng bcrypt và cấp JWT ngay sau khi hoàn tất) hoặc đăng ký nhanh qua Google OAuth; sau khi đăng ký thành công, hệ thống tự động đăng nhập và chuyển người học đến trang chủ. \\ \hline
    UC02 & Đăng nhập & Người học đăng nhập bằng email/mật khẩu hoặc tài khoản Google; hệ thống xác thực thông tin và cấp JWT access token để duy trì phiên làm việc. \\ \hline
    UC03 & Học từ vựng Foundation & Người học chọn cuốn sách và unit trong phân hệ Foundation Vocabulary (loạt \textit{4000 Essential English Words}, 6 cuốn $\times$ 30 unit), xem danh sách từ kèm định nghĩa và ví dụ, làm hai bài tập củng cố (20 câu mỗi bài), hoàn thành bài đọc hiểu cuối unit (5 câu hỏi); hệ thống lưu tiến độ từng unit và hỗ trợ bôi đen từ để gửi sang Vocab Lab bằng Lexon AI. \\ \hline
    UC04 & Học ngữ pháp Foundation & Người học truy cập bài học ngữ pháp theo cấp độ (Elementary/Intermediate/Advanced) trong phân hệ Foundation Grammar, đọc phần lý thuyết rồi làm bài tập thực hành (trắc nghiệm, điền vào chỗ trống, sắp xếp câu); hệ thống tự động chấm điểm và lưu tiến độ từng unit. \\ \hline
    UC05 & Luyện phát âm Foundation & Người học chọn âm IPA trong bộ 44 âm vị, xem mô tả khẩu hình và nghe audio chuẩn, sau đó ghi âm giọng mình để hệ thống phân tích và trả về điểm chính xác; FREE giới hạn 5 lượt ghi âm mỗi ngày. \\ \hline
    UC06 & Học IELTS Basic --- Listening & Người học học lý thuyết các dạng câu hỏi Listening IELTS (phần Lessons kèm quiz củng cố) và luyện bài tập thực hành chuẩn IELTS (phần Exercises); hệ thống tự động chấm điểm và lưu tiến độ; FREE giới hạn 3 Lessons và 2 Exercises mỗi kỹ năng. \\ \hline
    UC07 & Học IELTS Basic --- Reading & Người học học lý thuyết dạng câu hỏi Reading và luyện bài tập thực hành với đoạn văn học thuật kèm các dạng câu hỏi chuẩn IELTS (True/False/Not Given, Matching Headings, điền từ, trắc nghiệm); hệ thống chấm điểm và lưu kết quả. \\ \hline
    UC08 & Học IELTS Basic --- Writing & Người học học lý thuyết cách viết Task~1 và Task~2 IELTS, luyện bài tập viết cơ bản; hệ thống chấm điểm tự động và lưu trạng thái hoàn thành. \\ \hline
    UC09 & Học IELTS Basic --- Speaking & Người học học lý thuyết về từng Part Speaking IELTS và luyện bài tập ghi âm câu trả lời; hệ thống lưu bản ghi âm và đánh dấu hoàn thành bài. \\ \hline
    UC10 & Luyện IELTS Advanced --- Listening & Người học (yêu cầu Premium/Pro) chọn Part Listening nâng cao (Part~1--4) mô phỏng sát đề IELTS thực tế, nghe audio và trả lời các dạng câu hỏi chuẩn; hệ thống chấm điểm tự động và lưu kết quả chi tiết. \\ \hline
    UC11 & Luyện IELTS Advanced --- Reading & Người học (yêu cầu Premium/Pro) chọn Part Reading (Part~1--3) với đoạn văn học thuật dài, trả lời các dạng câu hỏi chuẩn IELTS trên giao diện hai cột; hệ thống chấm điểm, lưu kết quả và hiển thị giải thích đáp án. \\ \hline
    UC12 & Luyện IELTS Advanced --- Writing & Người học (yêu cầu Premium/Pro) chọn đề Writing (Task~1 hoặc Task~2), viết bài và nộp; hệ thống gửi bài vào hàng đợi RabbitMQ, backend-ai dùng Gemini chấm theo 4 tiêu chí IELTS Writing, trả về band điểm và nhận xét chi tiết qua HTTP callback; Premium tối đa 10 lượt/tháng, Pro không giới hạn. \\ \hline
    UC13 & Luyện IELTS Advanced --- Speaking & Người học (yêu cầu Premium/Pro) luyện Speaking theo Part~1/2/3, ghi âm câu trả lời và nộp; hệ thống gửi bài lên Gemini qua RabbitMQ để chấm theo 4 tiêu chí IELTS Speaking và trả về band điểm kèm nhận xét; có thể xem bài chia sẻ từ cộng đồng và bài đã nộp của mình. \\ \hline
    UC14 & Thi thử Listening (Mock Test) & Người học thi thử phần Listening toàn phần theo chuẩn IELTS (4 phần, 40 câu, đồng hồ đếm ngược theo thời gian quy định); hệ thống chấm điểm tự động và trả về band điểm ước tính sau khi nộp. \\ \hline
    UC15 & Thi thử Reading (Mock Test) & Người học thi thử phần Reading toàn phần (3 đoạn văn, 40 câu, 60 phút); hệ thống chấm điểm tự động và trả về band điểm ước tính. \\ \hline
    UC16 & Thi thử Writing (Mock Test) & Người học thi thử phần Writing toàn phần (Task~1 + Task~2, 60 phút); bài được gửi vào RabbitMQ để Gemini chấm bất đồng bộ, trả về band điểm và nhận xét theo từng tiêu chí. \\ \hline
    UC17 & Thi thử Speaking (Mock Test) & Người học tiến hành thi thử Speaking toàn phần mô phỏng 3 Part với AI, ghi âm câu trả lời từng câu theo đồng hồ đếm ngược chuẩn IELTS; bài được gửi vào RabbitMQ để Gemini chấm bất đồng bộ, trả về band điểm và nhận xét theo bốn tiêu chí IELTS Speaking. \\ \hline
    UC18 & Xem lịch sử bài thi (Test History) & Người học xem danh sách toàn bộ phiên thi đã thực hiện trong IELTS Intensive, lọc theo kỹ năng hoặc thời gian, xem chi tiết band điểm và nhận xét từng bài, xóa phiên thi cũ nếu cần; FREE chỉ lưu tối đa 3 phiên gần nhất. \\ \hline
    UC19 & Xây dựng lộ trình học tập (Roadmap) & Người học hoàn thành onboarding (chọn band mục tiêu, mức cam kết học hàng ngày, ngày thi dự kiến), tùy chọn làm bài kiểm tra chẩn đoán 3 kỹ năng; hệ thống sinh lộ trình học cá nhân hóa theo tuần và tự động đánh dấu bước hoàn thành khi đạt $\geq 90\%$. \\ \hline
    UC20 & Xem thống kê tổng quan (Statistics) & Người học xem trang Overview trong Statistics: band điểm ước tính từng kỹ năng, mục tiêu học hàng ngày, heatmap hoạt động trong năm và đếm ngược ngày thi. \\ \hline
    UC21 & Xem thống kê theo cấp độ & Người học xem thống kê chi tiết phân theo bốn tab: Foundation (tiến độ Vocab/Grammar/Pronunciation), Basic (lesson và exercise theo kỹ năng), Advanced (phiên luyện theo Part) và Intensive (lịch sử band điểm các phiên thi thử). \\ \hline
    UC22 & Quy đổi và tra cứu điểm IELTS (Calculator) & Người học nhập số câu đúng để tra band Listening hoặc Reading theo bảng chuẩn IELTS; xem descriptors tiêu chí band Writing/Speaking; nhập band bốn kỹ năng để tính band tổng theo quy tắc làm tròn chính thức. \\ \hline
    UC23 & Luyện Shadowing & Người học chọn bài video từ thư viện Shadowing (lọc theo danh mục, trạng thái cá nhân), nghe và nói đuổi theo từng câu transcript đồng bộ với video YouTube; hệ thống lưu tiến độ từng câu và cập nhật chuỗi học liên tiếp (streak). \\ \hline
    UC24 & Luyện Dictation & Người học chọn bài Dictation, chọn cấp độ khó (Beginner/Intermediate/Advanced/Expert), nghe audio từng câu và điền từ bị ẩn vào ô trống với hỗ trợ gợi ý tiệm tiến; kiểm tra đáp án từng câu và lưu tiến độ hoàn thành. \\ \hline
    UC25 & Tạo bài Shadowing cá nhân & Người học (yêu cầu Premium/Pro) nhập URL YouTube để tạo bài Shadowing riêng; backend-ai dùng faster-whisper phiên âm và phân tách transcript, cập nhật trạng thái từ PROCESSING sang READY; giao diện polling mỗi 5 giây và mở bài khi sẵn sàng. \\ \hline
    UC26 & Tạo bài Dictation cá nhân & Người học (yêu cầu Premium/Pro) nhập URL YouTube để tạo bài Dictation riêng; backend-ai xử lý audio và sinh transcript có mốc thời gian qua faster-whisper, cập nhật trạng thái READY; người học luyện tập giống thư viện chuẩn. \\ \hline
    UC27 & Tạo loại thẻ flashcard (Card Type) & Người học tạo Card Type mới trong Vocab Lab: đặt tên, định nghĩa danh sách Fields (kiểu Text hoặc Media) và Templates (mẫu hiển thị hai mặt thẻ); lưu để sử dụng khi tạo flashcard mới. \\ \hline
    UC28 & Thêm từ vựng bằng Lexon AI & Người học bôi đen từ hoặc cụm từ trong nội dung bài học bất kỳ; Lexon AI tự động sinh nội dung flashcard theo cấu trúc Card Type đã chọn và điền sẵn vào tab Add của Vocab Lab; người học xem lại, chỉnh sửa nếu cần và lưu thẻ mà không rời khỏi trang đang học; Premium tối đa 50 lượt/tháng. \\ \hline
    UC29 & Quản lý bộ thẻ (Decks) & Người học tạo, đổi tên, xóa, nhập (từ file \texttt{.lexon}) và xuất bộ thẻ flashcard; mỗi deck hiển thị số thẻ New/Learning/Due; bắt đầu phiên ôn tập FSRS; FREE giới hạn 3 deck và 50 thẻ/deck. \\ \hline
    UC30 & Nhập bộ thẻ từ Marketplace & Người học (yêu cầu Premium/Pro) tìm kiếm và nhập bộ thẻ công khai từ Community Marketplace theo danh mục (IELTS, TOEFL, Academic...) về thư viện cá nhân để ôn luyện. \\ \hline
    UC31 & Xuất bản bộ thẻ lên Marketplace & Người học (yêu cầu Premium/Pro) xuất bản bộ thẻ của mình lên Community Marketplace để chia sẻ với cộng đồng kèm tags, mô tả và danh mục chủ đề. \\ \hline
    UC32 & Xem thống kê Vocab Lab (Stats) & Người học xem biểu đồ phân bố thẻ theo trạng thái SRS (New/Learning/Review/Relearning), lịch sử review theo ngày, dự báo số thẻ đến hạn trong 30 ngày tới và hoạt động theo khung giờ trong ngày. \\ \hline
    UC33 & Xem và quản lý thông báo & Người học mở danh sách thông báo theo thứ tự mới nhất (kết quả chấm Writing/Speaking từ AI, cập nhật streak, sự kiện gói thành viên, thông báo từ Admin); đánh dấu đã đọc từng mục hoặc tất cả; nhấn vào thông báo để điều hướng đến trang liên quan. \\ \hline
    UC34 & Xem và tương tác bài đăng Community & Người học xem feed bài đăng cộng đồng (lọc theo All Posts/Study Tips/Achievements), tương tác bằng Like, Bookmark, Comment và Reply. \\ \hline
    UC35 & Đăng bài trong Community & Người học tạo bài đăng kèm tiêu đề, nội dung, tag phân loại và tối đa 4 ảnh đính kèm (ảnh lưu trên Cloudinary); bài xuất hiện đầu feed ngay sau khi đăng và người học nhận +5~XP; lần đăng đầu tiên nhận thêm huy hiệu First Post. \\ \hline
    UC36 & Xem bảng xếp hạng (Leaderboard) & Người học xem hai bảng xếp hạng top~10 trong Community: bảng XP tuần này và bảng Streak; dòng của người học đang đăng nhập được làm nổi bật để dễ nhận vị trí. \\ \hline
    UC37 & Đăng ký gói thành viên & Người học xem chi tiết và so sánh ba gói (FREE/PREMIUM/PRO), đăng ký và thanh toán qua cổng VNPay (HMAC-SHA512, chế độ sandbox); gói thành viên kích hoạt ngay sau khi thanh toán thành công; người dùng mới được dùng thử Premium 7 ngày miễn phí. \\ \hline
    UC38 & Tạo bài học Shadowing (Admin) & Người quản lý tạo bài Shadowing mới: nhập tiêu đề, ID video YouTube, thư mục, danh mục chủ đề, thời lượng và danh sách câu transcript kèm mốc thời gian bắt đầu/kết thúc (ms); xuất bản để bài hiển thị ngay trong thư viện người học. \\ \hline
    UC39 & Tạo bài học Dictation (Admin) & Người quản lý tạo bài Dictation mới: nhập ID video YouTube, cung cấp transcript chuẩn xác đến từng từ làm ground-truth để hệ thống so khớp đáp án; xuất bản cho người học. \\ \hline
    UC40 & Chỉnh sửa và xóa bài học (Admin) & Người quản lý chỉnh sửa nội dung (tiêu đề, transcript, chủ đề...) hoặc xóa bài học Shadowing và Dictation đã tạo; thay đổi có hiệu lực ngay lập tức với người học. \\ \hline
    UC41 & Quản lý gói đăng ký thành viên (Admin) & Người quản lý xem tổng quan số lượng tài khoản theo gói và doanh thu; xem danh sách chi tiết trạng thái đăng ký của từng học viên; cấp gói thành viên thủ công cho tài khoản cụ thể với thời hạn tùy chỉnh. \\ \hline
    UC42 & Liên kết và theo dõi học viên (Student/Teacher) & Người học đóng vai học viên nhập Teacher ID để liên kết với giáo viên; đóng vai giáo viên chia sẻ User ID của mình và xem tiến độ IELTS (band điểm, lịch sử thi, hoạt động học) của từng học viên đã liên kết. \\ \hline
    UC43 & Chat với trợ lý AI (mobile) & Người học trên ứng dụng di động mở giao diện chat để trò chuyện tự do với trợ lý Lexon AI về các chủ đề học tiếng Anh và IELTS; lịch sử hội thoại được duy trì trong suốt phiên làm việc. \\ \hline
    \caption{Danh sách các use case}
    \label{tab:uc_list}
\end{longtable}

\section{Đặc tả các yêu cầu chức năng}

Phần này trình bày đặc tả chi tiết cho 11 use case được lựa chọn từ tổng cộng 43 use case của hệ thống. Tiêu chí lựa chọn tập trung vào các use case có (i) mức độ phức tạp kỹ thuật cao --- ví dụ máy trạng thái dài như UC17 (Thi thử Speaking) hoặc pipeline AI bất đồng bộ như UC12 (Luyện Writing nâng cao); (ii) tính đặc trưng nghiệp vụ rõ nét --- ví dụ tra cứu từ vựng bằng AI (UC28), tạo loại thẻ tùy chỉnh (UC27), xây dựng lộ trình học cá nhân hóa (UC19); hoặc (iii) tần suất sử dụng cao trong vòng đời tài khoản --- ví dụ đăng ký (UC01), đăng nhập (UC02), học từ vựng (UC03). Mỗi use case bao gồm bảng mô tả luồng sự kiện (chính và thay thế), sơ đồ activity, sơ đồ tuần tự và mô tả kỹ thuật chi tiết.

\subsection{Đăng ký tài khoản (UC01)}
Bảng~\ref{tab:uc01_spec} đặc tả luồng sự kiện; hình~\ref{fig:activity-uc01} thể hiện sơ đồ activity; hình~\ref{fig:sequence-uc01} và hình~\ref{fig:sequence-uc01-gg} thể hiện sơ đồ tuần tự cho hai luồng đăng ký bằng email và Google OAuth.
\subsubsection{Mô tả use case}
\begin{table}[H]
    \centering
    \renewcommand{\arraystretch}{1.3}
    \begin{tabularx}{\linewidth}{|X|X|}
        \hline
        \multicolumn{2}{|l|}{\textbf{Đặc tả use-case}} \\ \hline
        \multicolumn{2}{|l|}{\textbf{Mã use-case:} UC01} \\ \hline
        \multicolumn{2}{|l|}{\textbf{Actor:} Người học} \\ \hline
        \multicolumn{2}{|p{\dimexpr\linewidth-2\tabcolsep-2\arrayrulewidth}|}{\textbf{Tiền điều kiện (Pre-condition):} Người dùng chưa có tài khoản và đang ở màn hình Đăng ký.} \\ \hline
        \multicolumn{2}{|p{\dimexpr\linewidth-2\tabcolsep-2\arrayrulewidth}|}{\textbf{Hậu điều kiện (Post-condition):} Tài khoản được tạo thành công; hệ thống tự động đăng nhập, cấp JWT access token và chuyển người học đến trang Home (áp dụng cho cả hai luồng Email và Google).} \\ \hline
        \multicolumn{2}{|l|}{\textbf{Luồng sự kiện chính (Basic flow)}} \\ \hline
        \textbf{Actor} & \textbf{Hệ thống} \\ \hline
        1a. Nhập Họ và tên, Email, Mật khẩu (tối thiểu 6 ký tự) và Xác nhận mật khẩu vào form, sau đó nhấn ``Create Account''. & \\ \hline
        & 2a. Kiểm tra tất cả các trường đã điền và mật khẩu xác nhận khớp nhau. \\ \hline
        & 3a. Kiểm tra định dạng email và xác nhận email chưa tồn tại trong cơ sở dữ liệu. \\ \hline
        & 4a. Băm mật khẩu, tạo tài khoản với vai trò người học và khởi tạo bộ thẻ từ vựng mặc định. \\ \hline
        & 5a. Đăng ký thành công; hệ thống tự động đăng nhập, cấp JWT access token và chuyển người học đến trang Home. Kết thúc use-case. \\ \hline
        1b. Nhấn ``Continue with Google'' và chọn tài khoản Google. & \\ \hline
        & 2b. Xác thực thông tin Google OAuth; tạo hoặc liên kết tài khoản từ thông tin Google (email, họ tên, ảnh đại diện), tạo tài khoản với vai trò người học và khởi tạo bộ thẻ mặc định. \\ \hline
        & 3b. Cấp JWT access token và chuyển đến trang Home. Kết thúc use-case. \\ \hline
        \multicolumn{2}{|l|}{\textbf{Luồng sự kiện thay thế (Alternative flow)}} \\ \hline
        & 2a.1. Một hoặc nhiều trường bỏ trống hoặc mật khẩu không khớp: thông báo lỗi tương ứng, không gửi yêu cầu lên máy chủ. Quay lại bước 1a. \\ \hline
        & 3a.1. Email đã tồn tại trong hệ thống: máy chủ trả về thông báo ``Email đã tồn tại''. Quay lại bước 1a. \\ \hline
        1b.1. Hủy đăng ký hoặc xác thực Google thất bại. & \\ \hline
        & 1b.2. Hiển thị thông báo lỗi, không tạo tài khoản. \\ \hline
    \end{tabularx}
    \caption{Đặc tả chức năng Đăng ký tài khoản (UC01)}
    \label{tab:uc01_spec}
\end{table}

\subsubsection{Sơ đồ activity}
\begin{figure}[H]
    \centering
    \includegraphics[width=\linewidth]{activity-diagram/activity-uc01.png}
    \caption{Sơ đồ activity UC01 -- Đăng ký tài khoản}
    \label{fig:activity-uc01}
\end{figure}

\subsubsection{Sơ đồ tuần tự}
\begin{figure}[H]
    \centering
    \includegraphics[width=\linewidth]{sequence-diagram/sequence-uc01.png}
    \caption{Sơ đồ tuần tự UC01 -- Đăng ký bằng email}
    \label{fig:sequence-uc01}
\end{figure}

\begin{figure}[H]
    \centering
    \includegraphics[width=\linewidth]{sequence-diagram/sequence-uc01-gg.png}
    \caption{Sơ đồ tuần tự UC01 -- Đăng ký bằng Google OAuth}
    \label{fig:sequence-uc01-gg}
\end{figure}

\subsubsection{Mô tả kỹ thuật chi tiết}

Hệ thống cung cấp hai phương thức đăng ký độc lập: đăng ký bằng email/mật khẩu và đăng ký qua Google OAuth.

Đối với luồng email, điểm cốt lõi là cơ chế bảo mật mật khẩu. Sau khi người học gửi thông tin đăng ký, máy chủ kiểm tra tính hợp lệ của email và xác nhận email chưa tồn tại trong cơ sở dữ liệu. Nếu hợp lệ, mật khẩu được băm một chiều với độ khó được hiệu chỉnh để cân bằng giữa bảo mật và hiệu năng xử lý. Tiếp theo, hệ thống tạo tài khoản mới với vai trò người học và tự động khởi tạo một bộ thẻ từ vựng mặc định phục vụ Vocab Lab. Nếu email đã tồn tại, hệ thống trả về thông báo lỗi tương ứng. Sau khi đăng ký thành công qua luồng email, hệ thống tự động đăng nhập và cấp JWT access token, chuyển người học đến trang Home mà không cần thực hiện thêm bước đăng nhập riêng.

\begin{sloppypar}
Đối với luồng Google OAuth, người học xác thực thông qua tài khoản Google trên thiết bị. Máy chủ xác minh thông tin xác thực từ Google, hỗ trợ đồng thời nhiều loại ứng dụng (web, iOS, Android) để tương thích đa nền tảng. Sau khi xác thực thành công, hệ thống áp dụng chiến lược tìm kiếm hoặc tạo mới tài khoản dựa trên thông tin định danh Google: nếu tài khoản chưa tồn tại thì tạo tài khoản mới và khởi tạo bộ thẻ mặc định. Khác với luồng email, luồng Google trả về JWT access token ngay lập tức, cho phép người dùng vào ứng dụng mà không cần thêm bước đăng nhập.
\end{sloppypar}

\subsection{Đăng nhập (UC02)}
Bảng~\ref{tab:uc02_spec} đặc tả luồng sự kiện; hình~\ref{fig:activity-uc02} thể hiện sơ đồ activity; hình~\ref{fig:sequence-uc02} và hình~\ref{fig:sequence-uc02-gg} thể hiện sơ đồ tuần tự cho hai luồng đăng nhập bằng email và Google OAuth.
\subsubsection{Mô tả use case}
\begin{table}[H]
    \centering
    \renewcommand{\arraystretch}{1.3}
    \begin{tabularx}{\linewidth}{|X|X|}
        \hline
        \multicolumn{2}{|l|}{\textbf{Đặc tả use-case}} \\ \hline
        \multicolumn{2}{|l|}{\textbf{Mã use-case:} UC02} \\ \hline
        \multicolumn{2}{|l|}{\textbf{Actor:} Người học} \\ \hline
        \multicolumn{2}{|p{\dimexpr\linewidth-2\tabcolsep-2\arrayrulewidth}|}{\textbf{Tiền điều kiện (Pre-condition):} Người học đã có tài khoản đăng ký trong hệ thống và đang ở màn hình Đăng nhập.} \\ \hline
        \multicolumn{2}{|p{\dimexpr\linewidth-2\tabcolsep-2\arrayrulewidth}|}{\textbf{Hậu điều kiện (Post-condition):} Người học truy cập thành công vào Dashboard; hệ thống cấp JWT.} \\ \hline
        \multicolumn{2}{|l|}{\textbf{Luồng sự kiện chính (Basic flow)}} \\ \hline
        \textbf{Actor} & \textbf{Hệ thống} \\ \hline
        1a. Nhập địa chỉ email và mật khẩu, sau đó nhấn ``Sign in''. & \\ \hline
        & 2a. Truy vấn cơ sở dữ liệu bằng email và xác minh mật khẩu nhập vào khớp với thông tin đã lưu. \\ \hline
        & 3a. Cấp JWT access token và chuyển hướng người học đến màn hình Dashboard. Kết thúc use-case. \\ \hline
        1b. Nhấn ``Đăng nhập với Google'' và chọn tài khoản Google. & \\ \hline
        & 2b. Xác thực thông tin Google OAuth; tìm hoặc tạo tài khoản tương ứng trong cơ sở dữ liệu. \\ \hline
        & 3b. Cấp JWT access token và chuyển hướng người học đến màn hình Dashboard. Kết thúc use-case. \\ \hline
        \multicolumn{2}{|l|}{\textbf{Luồng sự kiện thay thế (Alternative flow)}} \\ \hline
        & 2a.1. Email không tồn tại hoặc mật khẩu sai: hiển thị thông báo ``Email hoặc mật khẩu không đúng'' và dừng luồng. Quay lại bước 1a. \\ \hline
        1b.1. Hủy đăng nhập hoặc xác thực Google thất bại. & \\ \hline
        & 1b.2. Hiển thị thông báo lỗi, không cấp JWT. \\ \hline
    \end{tabularx}
    \caption{Đặc tả chức năng Đăng nhập (UC02)}
    \label{tab:uc02_spec}
\end{table}

\subsubsection{Sơ đồ activity}
\begin{figure}[H]
    \centering
    \includegraphics[width=\linewidth]{activity-diagram/activity-uc02.png}
    \caption{Sơ đồ activity UC02 -- Đăng nhập}
    \label{fig:activity-uc02}
\end{figure}

\subsubsection{Sơ đồ tuần tự}
\begin{figure}[H]
    \centering
    \includegraphics[width=\linewidth]{sequence-diagram/sequence-uc02.png}
    \caption{Sơ đồ tuần tự UC02 -- Đăng nhập bằng email}
    \label{fig:sequence-uc02}
\end{figure}

\begin{figure}[H]
    \centering
    \includegraphics[width=\linewidth]{sequence-diagram/sequence-uc02-gg.png}
    \caption{Sơ đồ tuần tự UC02 -- Đăng nhập bằng Google OAuth}
    \label{fig:sequence-uc02-gg}
\end{figure}

\subsubsection{Mô tả kỹ thuật chi tiết}

Điểm kỹ thuật cốt lõi trong UC02 là quá trình xác thực mật khẩu. Hệ thống sử dụng cơ chế so sánh theo thời gian hằng định, tức là thời gian thực thi không phụ thuộc vào vị trí ký tự sai, giúp ngăn chặn các cuộc tấn công phân tích thời gian. Luồng xác thực email/mật khẩu được xử lý qua một chiến lược xác thực riêng: hệ thống truy vấn cơ sở dữ liệu theo email, so sánh mật khẩu với giá trị đã băm và đính kèm thông tin người dùng vào yêu cầu trước khi cấp token.

\begin{sloppypar}
JWT access token được ký bằng khóa bí mật và có hiệu lực 7 ngày; khi token hết hạn, người học cần đăng nhập lại để nhận token mới. Mọi endpoint yêu cầu xác thực đều được bảo vệ bởi cơ chế kiểm tra JWT để giải mã token và xác nhận tính hợp lệ. Ngoài luồng email/mật khẩu, hệ thống hỗ trợ đăng nhập Google OAuth: người học gửi thông tin xác thực Google, máy chủ xác minh tính hợp lệ từ phía Google, sau đó tìm kiếm hoặc tạo tài khoản tương ứng trong cơ sở dữ liệu và cấp JWT theo cùng quy trình.
\end{sloppypar}

% ============================================================
% UC03 – ADD CARD TYPE
% ============================================================
\subsection{Tạo loại thẻ học (UC27)}
Bảng~\ref{tab:uc27_spec} đặc tả luồng sự kiện; hình~\ref{fig:activity-uc27} thể hiện sơ đồ activity; hình~\ref{fig:sequence-uc27} thể hiện sơ đồ tuần tự.
\subsubsection{Mô tả use case}
\begin{table}[H]
    \centering
    \renewcommand{\arraystretch}{1.3}
    \begin{tabularx}{\linewidth}{|X|X|}
        \hline
        \multicolumn{2}{|l|}{\textbf{Đặc tả use-case}} \\ \hline
        \multicolumn{2}{|l|}{\textbf{Mã use-case:} UC27} \\ \hline
        \multicolumn{2}{|l|}{\textbf{Actor:} Người học} \\ \hline
        \multicolumn{2}{|p{\dimexpr\linewidth-2\tabcolsep-2\arrayrulewidth}|}{\textbf{Tiền điều kiện (Pre-condition):} Người học đã đăng nhập thành công và đang ở màn hình Vocab Lab.} \\ \hline
        \multicolumn{2}{|p{\dimexpr\linewidth-2\tabcolsep-2\arrayrulewidth}|}{\textbf{Hậu điều kiện (Post-condition):} Loại thẻ mới được tạo và lưu trong hệ thống; người học có thể sử dụng loại thẻ này để tạo flashcard.} \\ \hline
        \multicolumn{2}{|l|}{\textbf{Luồng sự kiện chính (Basic flow)}} \\ \hline
        \textbf{Actor} & \textbf{Hệ thống} \\ \hline
        1. Nhấn ``Add'', sau đó nhấn vào selector ``Type'' ở form thêm thẻ. & \\ \hline
        & 2. Hiển thị màn hình ``Choose Card Type'' với danh sách loại thẻ hiện có. \\ \hline
        3. Nhấn biểu tượng cài đặt. & \\ \hline
        & 4. Tải danh sách các loại thẻ của người học từ cơ sở dữ liệu và hiển thị màn hình ``Manage Card Types''. \\ \hline
        5. Nhấn nút dấu cộng (+). & \\ \hline
        & 6. Hiển thị form nhập tên inline với ô văn bản (placeholder: ``Card type name\ldots''). \\ \hline
        7. Nhập tên loại thẻ rồi nhấn nút "create". & \\ \hline
        & 8. Tạo loại thẻ mới trong cơ sở dữ liệu và tự động mở màn hình ``Edit Card Type'' với ba tab: \textbf{Fields}, \textbf{Templates}, \textbf{Preview}. \\ \hline
        9. Cấu hình loại thẻ qua ba tab: (a) tab \textit{Fields} --- thêm, sửa, xóa trường với tên, kiểu (\textit{Text} hoặc \textit{Media}) và mô tả; (b) tab \textit{Templates} --- chọn trường hiển thị mặt trước/mặt sau; định dạng từng trường (in đậm, nghiêng, gạch chân, cỡ chữ, căn lề, màu chữ); thiết lập kiểu thẻ tổng thể (màu nền, màu chữ); (c) tab \textit{Preview} --- xem trước thẻ trực tiếp. & \\ \hline
        10. Nhấn ``Save''. & \\ \hline
        & 11. Lưu toàn bộ cấu hình (trường, mẫu hiển thị, kiểu dáng) vào cơ sở dữ liệu; đóng màn hình và quay về danh sách loại thẻ. \\ \hline
        \multicolumn{2}{|l|}{\textbf{Luồng sự kiện thay thế (Alternative flow)}} \\ \hline
        & 7.1. Tên loại thẻ rỗng: hệ thống không cho phép xác nhận. Quay lại bước 7. \\ \hline
        10.1. Nhấn ``Cancel'' thay vì ``Save''. & \\ \hline
        & 10.2. Hệ thống hủy toàn bộ thay đổi chưa lưu và đóng màn hình. \\ \hline
    \end{tabularx}
    \caption{Đặc tả chức năng Tạo loại thẻ học (UC27)}
    \label{tab:uc27_spec}
\end{table}

\subsubsection{Sơ đồ activity}
\begin{figure}[H]
    \centering
    \includegraphics[width=\linewidth]{activity-diagram/activity-uc27.png}
    \caption{Sơ đồ activity UC27 -- Tạo loại thẻ học}
    \label{fig:activity-uc27}
\end{figure}

\subsubsection{Sơ đồ tuần tự}
\begin{figure}[H]
    \centering
    \includegraphics[width=\linewidth]{sequence-diagram/sequence-uc27.png}
    \caption{Sơ đồ tuần tự UC27 -- Tạo loại thẻ học}
    \label{fig:sequence-uc27}
\end{figure}

\subsubsection{Mô tả kỹ thuật chi tiết}

Chức năng tạo loại thẻ học (Card Type) là nền tảng của toàn bộ hệ thống Vocab Lab, cho phép người học tùy chỉnh cấu trúc flashcard theo nhu cầu học tập cá nhân. Luồng tạo loại thẻ được thiết kế theo hai giai đoạn tách biệt: khởi tạo tên và cấu hình chi tiết.

Ở giai đoạn khởi tạo, người học nhập tên loại thẻ qua form inline. Hệ thống kiểm tra tên không được rỗng trước khi cho phép xác nhận, sau đó tạo bản ghi loại thẻ mới trong cơ sở dữ liệu và lập tức chuyển sang giao diện chỉnh sửa chi tiết. Thiết kế này đảm bảo loại thẻ tồn tại trong cơ sở dữ liệu ngay từ đầu, giúp tránh mất dữ liệu nếu người học rời khỏi màn hình giữa chừng.

Giao diện chỉnh sửa chi tiết được tổ chức thành ba tab có chức năng độc lập nhưng liên kết chặt chẽ. Tab \textit{Fields} cho phép định nghĩa cấu trúc dữ liệu của thẻ: mỗi trường có tên, kiểu dữ liệu (văn bản hoặc tệp phương tiện) và mô tả ngắn. Tab \textit{Templates} ánh xạ các trường vừa định nghĩa lên bố cục hiển thị: người học chọn trường nào xuất hiện ở mặt trước, trường nào ở mặt sau, đồng thời cấu hình định dạng trình bày (font, màu sắc, căn lề) cho từng trường. Tab \textit{Preview} kết xuất thẻ theo đúng cấu hình hiện tại, cho phép người học xem trước kết quả mà không cần lưu.

Khi người học nhấn ``Save'', toàn bộ cấu hình --- bao gồm danh sách trường, mẫu hiển thị hai mặt và thiết lập kiểu dáng tổng thể --- được cập nhật vào cơ sở dữ liệu trong một thao tác duy nhất. Thiết kế lưu toàn bộ một lần thay vì lưu từng phần giúp đảm bảo tính nhất quán: loại thẻ không bao giờ ở trạng thái nửa chừng với cấu hình mâu thuẫn giữa các tab.

% ============================================================
% UC04 – LOOK UP AND ADD WORDS WITH AI
% ============================================================
\subsection{Tra cứu và thêm từ vựng bằng AI (UC28)}
Bảng~\ref{tab:uc28_spec} đặc tả luồng sự kiện; hình~\ref{fig:activity-uc28} thể hiện sơ đồ activity; hình~\ref{fig:sequence-uc28} thể hiện sơ đồ tuần tự.
\subsubsection{Mô tả use case}
\renewcommand{\arraystretch}{1.3}
\begin{longtable}{|p{0.46\linewidth}|p{0.46\linewidth}|}
    \caption{Đặc tả chức năng Tra cứu và thêm từ vựng bằng AI (UC28)} \label{tab:uc28_spec} \\
    \hline
    \multicolumn{2}{|l|}{\textbf{Đặc tả use-case}} \\ \hline
    \multicolumn{2}{|l|}{\textbf{Mã use-case:} UC28} \\ \hline
    \multicolumn{2}{|l|}{\textbf{Actor:} Người học} \\ \hline
    \multicolumn{2}{|p{0.88\linewidth}|}{\textbf{Tiền điều kiện (Pre-condition):} Người học đã đăng nhập thành công vào hệ thống web và đang xem nội dung bài học có văn bản. Vocab Lab của người học có ít nhất một loại thẻ (Card Type) đã được định nghĩa.} \\ \hline
    \multicolumn{2}{|p{0.88\linewidth}|}{\textbf{Hậu điều kiện (Post-condition):} Một Flashcard mới với nội dung do Lexon AI sinh ra được lưu vào Vocab Lab của người học, gắn với deck đã chọn.} \\ \hline
    \multicolumn{2}{|l|}{\textbf{Luồng sự kiện chính (Basic flow)}} \\ \hline
    \textbf{Actor} & \textbf{Hệ thống} \\ \hline
    \endfirsthead
    \multicolumn{2}{r}{\textit{(tiếp trang trước)}} \\
    \hline
    \textbf{Actor} & \textbf{Hệ thống} \\ \hline
    \endhead
    \hline
    \multicolumn{2}{r}{\textit{(tiếp trang sau)}} \\
    \endfoot
    \hline
    \endlastfoot
    1. Bôi đen từ hoặc cụm từ muốn tra cứu trong nội dung bài học. & \\ \hline
    & 2. Hệ thống phát hiện vùng văn bản được chọn, trích xuất từ cùng câu ngữ cảnh xung quanh, sau đó hiển thị menu ngữ cảnh nổi phía trên vùng chọn với hai tùy chọn: ``Add to Vocab'' và ``Ask AI''. \\ \hline
    3. Nhấn ``Ask AI''. & \\ \hline
    & 4. Mở cửa sổ Lexon AI kèm từ và câu ngữ cảnh; tải danh sách loại thẻ của người học từ cơ sở dữ liệu. \\ \hline
    & 5. Lexon AI hiển thị tin nhắn giới thiệu từ đã chọn kèm câu hỏi ``Please choose a format below to explain it:'' và các nút gợi ý ``Explain as [tên-loại-thẻ]'' tương ứng với từng loại thẻ. \\ \hline
    6. Nhấn nút ``Explain as [tên-loại-thẻ]''. & \\ \hline
    & 7. Hệ thống gửi yêu cầu đến AI với thông tin: tên loại thẻ, từ cần giải thích, câu ngữ cảnh và mô tả từng trường của loại thẻ; AI phân tích và trả về giải thích chi tiết theo đúng từng trường. Hệ thống ghi nhớ loại thẻ vừa chọn làm lựa chọn ưu tiên cho lần sau. \\ \hline
    & 8. Lexon AI hiển thị nội dung giải thích kèm hai nhóm nút gợi ý: ``Add to Vocab Lab'' và ``Explain as [tên-loại-thẻ-khác]'' cho các loại thẻ còn lại. \\ \hline
    9. Nhấn ``Add to Vocab Lab''. & \\ \hline
    & 10. Hệ thống gửi yêu cầu thứ hai đến AI để sinh dữ liệu thẻ theo cấu trúc JSON với các trường tương ứng với loại thẻ đã chọn; Lexon AI hiển thị thông báo ``Generating card data...''. \\ \hline
    & 11. Sau khi nhận dữ liệu từ AI, hệ thống mở form thêm thẻ nhanh và tự động điền sẵn nội dung theo đúng cấu trúc loại thẻ đã chọn. \\ \hline
    12. Xem lại, chỉnh sửa nội dung nếu cần, chọn deck đích và nhấn ``Add Card''. & \\ \hline
    & 13. Lưu flashcard vào cơ sở dữ liệu gắn với deck đã chọn; hiển thị thông báo ``Card added successfully!'' và làm mới form. Kết thúc use-case. \\ \hline
    \multicolumn{2}{|l|}{\textbf{Luồng sự kiện thay thế (Alternative flow)}} \\ \hline
    & 4.1. Người học chưa tạo loại thẻ nào: Lexon AI hiển thị tin nhắn chung về từ đã chọn để người học tự đặt câu hỏi tự do; không hiển thị nút gợi ý, thẻ không được tạo. Kết thúc use-case. \\ \hline
    & 5.1. Hệ thống phát hiện loại thẻ ưu tiên đã được ghi nhớ từ lần tra cứu trước: bỏ qua bước 5 và 6, tự động gọi AI giải thích ngay theo loại thẻ ưa thích đó. Tiếp tục từ bước 7. \\ \hline
    & 7.1. Lỗi kết nối hoặc phản hồi AI không hợp lệ ở bước giải thích: Lexon AI hiển thị thông báo ``Sorry, I failed to generate the explanation. Please try again.'' Quay lại bước 6. \\ \hline
    9.1. Nhấn ``Explain as [tên-loại-thẻ-khác]'' thay vì ``Add to Vocab Lab''. & \\ \hline
    & 9.2. Hệ thống gọi lại bước 7 với loại thẻ mới; Lexon AI giải thích theo cấu trúc loại thẻ vừa chọn và hiển thị lại bộ nút gợi ý tương ứng. \\ \hline
    & 10.1. AI trả về dữ liệu không hợp lệ hoặc lỗi kết nối ở bước sinh dữ liệu thẻ: Lexon AI hiển thị thông báo ``Sorry, I failed to generate the card fields. Please try again.'' Quay lại bước 9. \\ \hline
    12.1. Nhấn ``X'' để đóng form thêm thẻ nhanh. & \\ \hline
    & 12.2.1. Form đóng lại; flashcard không được lưu. Kết thúc use-case. \\ \hline
    & 12.2.2. Chưa chọn deck hoặc chưa có deck nào: hệ thống hiển thị thông báo ``Please select or create a deck first.''; người học cần tạo hoặc chọn deck rồi thực hiện lại. Quay lại bước 12. \\ \hline
\end{longtable}

\subsubsection{Sơ đồ activity}
\begin{figure}[H]
    \centering
    \includegraphics[width=\linewidth]{activity-diagram/activity-uc28.png}
    \caption{Sơ đồ activity UC28 -- Tra cứu và thêm từ vựng bằng AI}
    \label{fig:activity-uc28}
\end{figure}

\subsubsection{Sơ đồ tuần tự}
\begin{figure}[H]
    \centering
    \includegraphics[width=\linewidth]{sequence-diagram/sequence-uc28.png}
    \caption{Sơ đồ tuần tự UC28 -- Tra cứu và thêm từ vựng bằng AI}
    \label{fig:sequence-uc28}
\end{figure}

\subsubsection{Mô tả kỹ thuật chi tiết}

Tính năng tra cứu và thêm từ vựng bằng AI (Lexon AI) được thiết kế theo mô hình ba giai đoạn: phát hiện vùng chọn, sinh giải thích và điền sẵn thẻ.

Ở giai đoạn phát hiện vùng chọn, hệ thống lắng nghe sự kiện chọn văn bản trên toàn bộ trang nội dung bài học. Khi phát hiện vùng chọn không rỗng, hệ thống trích xuất từ được chọn cùng câu ngữ cảnh xung quanh (khoảng 80 ký tự) và hiển thị menu ngữ cảnh nổi ngay phía trên vùng chọn. Hai tùy chọn được cung cấp: ``Add to Vocab'' (thêm trực tiếp) và ``Ask AI'' (tra cứu bằng AI).

\begin{sloppypar}
Khi người học chọn ``Ask AI'', hệ thống mở cửa sổ trợ lý Lexon AI và đồng thời tải danh sách loại thẻ của người học từ cơ sở dữ liệu. Nếu người học đã dùng một loại thẻ cụ thể trong lần tra cứu trước, hệ thống ghi nhớ lựa chọn đó và tự động áp dụng cho lần này mà không cần hỏi lại. Khi người học xác nhận loại thẻ muốn dùng, hệ thống xây dựng yêu cầu gửi đến AI với đầy đủ ngữ cảnh: từ cần giải thích, câu ngữ cảnh xung quanh và mô tả cấu trúc từng trường của loại thẻ được chọn. Thiết kế prompt engineering chi tiết này giúp AI trả về giải thích đúng định dạng cấu trúc loại thẻ, thay vì phản hồi tự do khó xử lý.
\end{sloppypar}

Ở giai đoạn điền sẵn thẻ, sau khi người học xác nhận muốn thêm vào Vocab Lab, hệ thống gửi yêu cầu thứ hai đến AI để sinh dữ liệu theo định dạng có cấu trúc (JSON). Dữ liệu nhận được được dùng để điền sẵn vào form thêm thẻ nhanh, giúp người học xem lại và chỉnh sửa trước khi lưu chính thức. Thiết kế tách biệt bước giải thích và bước sinh dữ liệu thẻ cho phép người học đọc hiểu nội dung trước khi quyết định lưu, đồng thời xử lý độc lập lỗi ở từng bước mà không ảnh hưởng đến nhau.

% ============================================================
% UC05 – STUDY IELTS FOUNDATION VOCABULARY
% ============================================================
\subsection{Học từ vựng nền tảng IELTS (UC03)}
Bảng~\ref{tab:uc03_spec} đặc tả luồng sự kiện; hình~\ref{fig:activity-uc03} thể hiện sơ đồ activity; hình~\ref{fig:sequence-uc03} thể hiện sơ đồ tuần tự.
\subsubsection{Mô tả use case}
\renewcommand{\arraystretch}{1.3}
\begin{longtable}{|p{0.46\linewidth}|p{0.46\linewidth}|}
    \caption{Đặc tả chức năng Học từ vựng nền tảng IELTS (UC03)} \label{tab:uc03_spec} \\
    \hline
    \multicolumn{2}{|l|}{\textbf{Đặc tả use-case}} \\ \hline
    \multicolumn{2}{|l|}{\textbf{Mã use-case:} UC03} \\ \hline
    \multicolumn{2}{|l|}{\textbf{Actor:} Người học} \\ \hline
    \multicolumn{2}{|p{0.88\linewidth}|}{\textbf{Tiền điều kiện (Pre-condition):} Người học đã đăng nhập thành công và đang ở trang ielts/vocabulary.} \\ \hline
    \multicolumn{2}{|p{0.88\linewidth}|}{\textbf{Hậu điều kiện (Post-condition):} Unit được đánh dấu hoàn thành và lưu vào cơ sở dữ liệu; người học nhận 15~XP từ hệ thống gamification; các từ đã được đánh giá SRS được lưu vào Vocab Lab của người học.} \\ \hline
    \multicolumn{2}{|l|}{\textbf{Luồng sự kiện chính (Basic flow)}} \\ \hline
    \textbf{Actor} & \textbf{Hệ thống} \\ \hline
    \endfirsthead
    \multicolumn{2}{r}{\textit{(tiếp trang trước)}} \\
    \hline
    \textbf{Actor} & \textbf{Hệ thống} \\ \hline
    \endhead
    \hline
    \multicolumn{2}{r}{\textit{(tiếp trang sau)}} \\
    \endfoot
    \hline
    \endlastfoot
    1. Chọn một trong sáu đầu sách từ vựng trong danh sách, nhấn ``Start Learning''. & \\ \hline
    & 2. Tải thông tin sách và tiến độ học của người học từ cơ sở dữ liệu; hiển thị màn hình danh sách unit, mỗi unit kèm thanh tiến độ số từ đã học trên tổng số từ và trạng thái đọc hiểu. \\ \hline
    3. Chọn một unit bất kỳ trong danh sách. & \\ \hline
    & 4. Tải dữ liệu unit (được cache để tối ưu hiệu năng); hiển thị tab Lesson với Flashcard đầu tiên --- mặt trước gồm từ vựng, phiên âm IPA và hình ảnh minh họa. \\ \hline
    5. Nhấn vào thẻ Flashcard. & \\ \hline
    & 6. Thực hiện hoạt ảnh lật 3D; mặt sau thẻ hiển thị nghĩa của từ, câu ví dụ minh họa, bốn nút đánh giá SRS (``Again'', ``Hard'', ``Good'', ``Easy'') và nút ``Already know --- skip''. \\ \hline
    7a. Chọn một trong bốn mức đánh giá SRS: ``Again'' (chưa nhớ), ``Hard'' (nhớ khó), ``Good'' (nhớ đúng) hoặc ``Easy'' (nhớ rất tốt). & \\ \hline
    & 8a. Tạo Flashcard trong bộ thẻ mặc định của Vocab Lab với trạng thái SRS tương ứng với mức đánh giá; hiển thị thông báo ``Added to Vocab Lab''; cập nhật tiến độ học. Chuyển sang Flashcard tiếp theo. \\ \hline
    7b. Nhấn ``Already know --- skip''. & \\ \hline
    & 8b. Chuyển sang Flashcard tiếp theo mà không thêm từ vào Vocab Lab; cập nhật tiến độ học. \\ \hline
    \multicolumn{2}{|p{0.88\linewidth}|}{Lặp lại bước 5--8 cho đến khi tất cả các từ trong unit được xem qua.} \\ \hline
    & 10. Tự động mở khóa tab ``Reading Comprehension''; hiển thị đoạn văn đọc hiểu với các từ đã học được in đậm kèm tooltip, và bộ câu hỏi kiểm tra bên phải (dạng điền từ hoặc trắc nghiệm). \\ \hline
    11. Đọc đoạn văn, điền câu trả lời cho từng câu hỏi, nhấn ``Submit Answers''. & \\ \hline
    & 12. Đối chiếu từng đáp án với đáp án chuẩn và trả về kết quả gồm số câu đúng, tổng số câu và trạng thái từng câu. \\ \hline
    & 13a. Nếu chưa đạt điểm tuyệt đối: hiển thị kết quả, tô màu đỏ câu sai và màu xanh câu đúng; hiển thị nút ``Try Again''. Quay lại bước 11. \\ \hline
    & 13b. Nếu toàn bộ câu trả lời đúng: đánh dấu unit hoàn thành và lưu vào cơ sở dữ liệu; kích hoạt sự kiện gamification 15~XP; hiển thị thông báo hoàn thành và nút ``Continue Learning'' để quay về danh sách unit. Kết thúc use-case. \\ \hline
    \multicolumn{2}{|l|}{\textbf{Luồng sự kiện thay thế (Alternative flow)}} \\ \hline
    7a.1. Chọn mức đánh giá ``Again'' (chưa nhớ từ). & \\ \hline
    & 7a.2. Chèn từ trở lại cuối hàng đợi, tổng số từ cần xem tăng thêm 1; Flashcard sẽ xuất hiện lại trong phiên học hiện tại. Quay lại bước 5. \\ \hline
    & 10.1. Unit đã được hoàn thành từ phiên trước: tab ``Reading Comprehension'' được mở sẵn ngay khi vào unit; người học xem lại đoạn văn và kết quả mà không cần làm lại bài. \\ \hline
    & 12.1. Kết nối mạng gián đoạn khi nộp đáp án: hiển thị thông báo lỗi, giữ nguyên câu trả lời để người học thử lại. Quay lại bước 11. \\ \hline
\end{longtable}

\subsubsection{Sơ đồ activity}
\begin{figure}[H]
    \centering
    \includegraphics[width=\linewidth]{activity-diagram/activity-uc03.png}
    \caption{Sơ đồ activity UC03 -- Học từ vựng nền tảng IELTS}
    \label{fig:activity-uc03}
\end{figure}

\subsubsection{Sơ đồ tuần tự}
\begin{figure}[H]
    \centering
    \includegraphics[width=\linewidth]{sequence-diagram/sequence-uc03.png}
    \caption{Sơ đồ tuần tự UC03 -- Học từ vựng nền tảng IELTS}
    \label{fig:sequence-uc03}
\end{figure}

\subsubsection{Mô tả kỹ thuật chi tiết}

Hệ thống Foundation Vocabulary được thiết kế theo mô hình phân tầng dữ liệu: nội dung tĩnh (danh mục sách, unit, từ vựng và câu hỏi) được cache phía máy chủ để giảm tải truy vấn lặp lại, trong khi tiến độ cá nhân được ghi nhận vào cơ sở dữ liệu với ràng buộc duy nhất theo cặp người học -- unit, đảm bảo mỗi cặp chỉ có một bản ghi.

Khi người học đánh giá một từ bằng SRS, hệ thống thực hiện hai thao tác đồng thời: tạo Flashcard trong bộ thẻ mặc định của Vocab Lab với trạng thái SRS tương ứng với mức đánh giá được chọn, và cập nhật tiến độ học. Khi người học chọn ``Already know --- skip'', từ không được thêm vào Vocab Lab nhưng vẫn được tính vào tiến độ học.

Điểm kỹ thuật đặc biệt với lựa chọn ``Again'': từ được chèn lại vào cuối hàng đợi phía client mà không cần truy vấn lên máy chủ, đảm bảo từ xuất hiện lại trong phiên học hiện tại. Điều kiện mở khóa tab đọc hiểu được kiểm tra hoàn toàn phía client: khi số từ đã xem đạt đến tổng số từ của unit, tab ``Reading Comprehension'' tự động chuyển sang trạng thái hoạt động.

Trong phần đọc hiểu, hệ thống áp dụng so sánh không phân biệt hoa thường cho dạng điền từ và so sánh chính xác giá trị cho dạng trắc nghiệm. Hệ thống yêu cầu người học trả lời đúng toàn bộ câu hỏi (100\%) để đánh dấu unit hoàn thành. Sau khi đạt điều kiện, hệ thống cập nhật trạng thái hoàn thành vào cơ sở dữ liệu, trao 15~XP và kiểm tra điều kiện mở khóa các thành tích học tập tương ứng.

% ============================================================
% UC06 – STUDY IELTS BASIC LISTENING LESSONS
% ============================================================
\subsection{Học bài Listening IELTS cơ bản (UC06)}
Bảng~\ref{tab:uc06_spec} đặc tả luồng sự kiện; hình~\ref{fig:activity-uc06} thể hiện sơ đồ activity; hình~\ref{fig:sequence-uc06-lessons} và hình~\ref{fig:sequence-uc06-exercises} thể hiện sơ đồ tuần tự cho luồng học lý thuyết và luồng luyện bài tập nghe.
\subsubsection{Mô tả use case}
\renewcommand{\arraystretch}{1.3}
\begin{longtable}{|p{0.46\linewidth}|p{0.46\linewidth}|}
    \caption{Đặc tả chức năng Học bài Listening IELTS cơ bản (UC06)} \label{tab:uc06_spec} \\
    \hline
    \multicolumn{2}{|l|}{\textbf{Đặc tả use-case}} \\ \hline
    \multicolumn{2}{|l|}{\textbf{Mã use-case:} UC06} \\ \hline
    \multicolumn{2}{|l|}{\textbf{Actor:} Người học} \\ \hline
    \multicolumn{2}{|p{0.88\linewidth}|}{\textbf{Tiền điều kiện (Pre-condition):} Người học đã đăng nhập thành công.} \\ \hline
    \multicolumn{2}{|p{0.88\linewidth}|}{\textbf{Hậu điều kiện (Post-condition):} Bài học lý thuyết hoặc bài tập luyện nghe được đánh dấu hoàn thành và lưu vào cơ sở dữ liệu; người học nhận 10~XP khi hoàn thành bài học hoặc 15~XP khi hoàn thành bài tập với điểm tuyệt đối; chuỗi học liên tiếp được cập nhật.} \\ \hline
    \multicolumn{2}{|l|}{\textbf{Luồng sự kiện chính (Basic flow)}} \\ \hline
    \textbf{Actor} & \textbf{Hệ thống} \\ \hline
    \endfirsthead
    \multicolumn{2}{r}{\textit{(tiếp trang trước)}} \\
    \hline
    \textbf{Actor} & \textbf{Hệ thống} \\ \hline
    \endhead
    \hline
    \multicolumn{2}{r}{\textit{(tiếp trang sau)}} \\
    \endfoot
    \hline
    \endlastfoot
    1. Nhấn vào mục ``IELTS Basic'' trên thanh điều hướng bên trái (sidebar) của phân hệ IELTS. & \\ \hline
    & 2. Tải thống kê tiến độ kỹ năng; hiển thị trang thư viện kỹ năng gồm 4 thẻ: \textit{Listening}, \textit{Reading}, \textit{Writing}, \textit{Speaking}. Mỗi thẻ kỹ năng hiển thị hai phần: \textbf{Lessons} (bài học lý thuyết) và \textbf{Exercises} (bài tập luyện tập) kèm tiến độ hoàn thành tương ứng. \\ \hline
    3a. Nhấn vào phần ``Lessons'' của thẻ kỹ năng Listening. & \\ \hline
    & 4a. Tải danh sách bài học lý thuyết của kỹ năng Listening; hiển thị theo thứ tự chương, mỗi bài kèm tiêu đề và trạng thái hoàn thành. \\ \hline
    5a. Chọn một bài học từ danh sách. & \\ \hline
    & 6a. Tải nội dung bài học; hiển thị trang bài học với các khối nội dung lý thuyết (strategy, tips, traps, overview, section) được render từ Markdown và thanh mục lục điều hướng bên cạnh. Cuối trang là quiz trắc nghiệm tối đa 4 câu. \\ \hline
    7a. Đọc nội dung lý thuyết, trả lời các câu hỏi quiz và nhấn ``Submit''. & \\ \hline
    & 8a. Chấm điểm quiz phía client; ghi nhận bài học hoàn thành vào cơ sở dữ liệu; kích hoạt sự kiện gamification (10~XP) và cập nhật chuỗi học liên tiếp; cập nhật thanh tiến độ. Kết thúc nhánh A. \\ \hline
    3b. Nhấn vào phần ``Exercises'' của thẻ kỹ năng Listening. & \\ \hline
    & 4b. Tải danh sách bài tập luyện nghe; hiển thị kèm tiêu đề, loại dạng bài và trạng thái hoàn thành. \\ \hline
    5b. Chọn một bài tập từ danh sách. & \\ \hline
    & 6b. Tải dữ liệu bài tập; hiển thị giao diện bài tập gồm: trình phát âm thanh tùy chỉnh (hỗ trợ seek theo thời gian), khu vực câu hỏi được render theo dạng bài và thanh theo dõi tiến độ câu hỏi ở phía dưới màn hình. \\ \hline
    7b. Nghe audio, điền câu trả lời vào từng câu hỏi theo dạng bài tương ứng. & \\ \hline
    & 8b. Cập nhật trạng thái đã điền của từng câu hỏi trên thanh theo dõi theo thời gian thực, xử lý hoàn toàn phía client. \\ \hline
    9b. Nhấn ``Submit Answers'' sau khi hoàn thành tất cả câu hỏi. & \\ \hline
    & 10b. Tính điểm: so sánh không phân biệt hoa thường cho dạng điền từ, so sánh chính xác giá trị cho trắc nghiệm, hỗ trợ nhiều đáp án chấp nhận được cho Short Answer. Hiển thị kết quả số câu đúng trên tổng số câu. \\ \hline
    & 11b. Nếu tất cả câu trả lời đúng (điểm tuyệt đối): ghi nhận bài tập hoàn thành vào cơ sở dữ liệu; kích hoạt sự kiện gamification (15~XP) và cập nhật chuỗi học liên tiếp; mở khóa bảng Transcript bên phải với nội dung đồng bộ theo câu hỏi; hiển thị nút ``Next Step''. Kết thúc use-case. \\ \hline
    \multicolumn{2}{|l|}{\textbf{Luồng sự kiện thay thế (Alternative flow)}} \\ \hline
    & 10b.1. Điểm số chưa đạt tuyệt đối: hiển thị số câu đúng và nút ``Retry''; không ghi nhận hoàn thành; bảng Transcript không được mở khóa. Quay lại bước 7b. \\ \hline
    & 8a.1. Bài học đã được hoàn thành từ lần học trước: giao diện hiển thị trạng thái đã hoàn thành; người học xem lại lý thuyết mà không cần làm lại quiz. \\ \hline
    & 6b.1. Câu hỏi trong bài tập được gắn mốc thời gian audio: nút ``Listen from here'' xuất hiện cạnh câu hỏi tương ứng; nhấn nút này sẽ seek trình phát đến đúng mốc thời gian trong audio mà không phát lại từ đầu. \\ \hline
\end{longtable}

\subsubsection{Sơ đồ activity}
\begin{figure}[H]
    \centering
    \includegraphics[width=\linewidth]{activity-diagram/activity-uc06.png}
    \caption{Sơ đồ activity UC06 -- Học bài Listening IELTS cơ bản}
    \label{fig:activity-uc06}
\end{figure}

\subsubsection{Sơ đồ tuần tự}
\begin{figure}[H]
    \centering
    \includegraphics[width=\linewidth]{sequence-diagram/sequence-uc06-lessons.png}
    \caption{Sơ đồ tuần tự UC06 -- Học bài lý thuyết}
    \label{fig:sequence-uc06-lessons}
\end{figure}

\begin{figure}[H]
    \centering
    \includegraphics[width=\linewidth]{sequence-diagram/sequence-uc06-exercises.png}
    \caption{Sơ đồ tuần tự UC06 -- Luyện bài tập nghe}
    \label{fig:sequence-uc06-exercises}
\end{figure}

\subsubsection{Mô tả kỹ thuật chi tiết}

\begin{sloppypar}
Phân hệ IELTS Basic Listening tổ chức nội dung theo hai tầng tách biệt: bài học lý thuyết và bài tập luyện nghe, cả hai đều liên kết với kỹ năng Listening. Tiến độ của mỗi người học được ghi nhận riêng lẻ với ràng buộc duy nhất theo cặp người học -- mục học, đảm bảo mỗi cặp chỉ có đúng một bản ghi; thao tác ghi sử dụng cơ chế tạo mới hoặc cập nhật để tránh xung đột khi người học hoàn thành lại.
\end{sloppypar}

Nội dung bài học được lưu dưới dạng cấu trúc JSON gồm mảng các khối, trong đó mỗi khối xác định loại nội dung (tổng quan, chiến lược làm bài, mẹo, bẫy thường gặp, phần nội dung) và chuỗi Markdown tương ứng. Giao diện tự động xây dựng thanh mục lục điều hướng dựa trên danh sách khối. Quiz cuối bài được chấm điểm hoàn toàn phía client trước khi ghi nhận hoàn thành lên máy chủ.

\begin{sloppypar}
Cấu trúc câu hỏi của bài tập Listening được lưu dưới dạng mảng các nhóm câu hỏi, mỗi nhóm xác định loại dạng bài và danh sách câu hỏi tương ứng. Hệ thống hỗ trợ chín dạng bài chuẩn IELTS Listening: \textit{Form/Note Completion}, \textit{Table Completion}, \textit{Flow Chart Completion}, \textit{Short Answer}, \textit{Matching}, \textit{Map/Plan Labelling}, \textit{Diagram Labelling}, \textit{Multiple Choice (Single)} và \textit{Multiple Choice (Multiple)}. Mỗi dạng bài được render bởi component riêng biệt, giúp tách biệt logic hiển thị và dễ mở rộng.
\end{sloppypar}

Điểm kỹ thuật đặc biệt là tính năng seek theo thời gian: mỗi câu hỏi có thể được gắn mốc thời gian audio, cho phép trình phát âm thanh tùy chỉnh hiển thị nút ``Listen from here''; khi nhấn, trình phát seek đến đúng mốc thời gian đó trong audio thay vì phát lại từ đầu, giúp người học tập trung vào đoạn nghe liên quan đến câu hỏi cụ thể.

\begin{sloppypar}
Logic chấm điểm được thực hiện hoàn toàn phía client: so sánh không phân biệt hoa thường và loại bỏ khoảng trắng thừa cho dạng điền từ; so sánh mảng lựa chọn cho dạng trắc nghiệm nhiều đáp án; hỗ trợ danh sách đáp án chấp nhận được cho Short Answer có nhiều phương án đúng. Chỉ khi đạt điểm tuyệt đối (toàn bộ câu hỏi đúng), hệ thống mới ghi nhận hoàn thành và trao 15~XP --- thiết kế này đảm bảo người học thực sự nắm vững nội dung trước khi được tính là hoàn thành, phù hợp với yêu cầu luyện thi nghiêm túc.
\end{sloppypar}

Sau khi đạt điểm tuyệt đối, bảng Transcript bên phải được mở khóa: nội dung transcript được hiển thị đồng bộ với số câu hỏi tương ứng, cho phép người học đối chiếu câu trả lời với vị trí thông tin trong bài nghe. Cả hai loại hoàn thành đều kích hoạt cơ chế cập nhật chuỗi học liên tiếp và kiểm tra điều kiện mở khóa thành tích học tập.

% ============================================================
% UC07 – PRACTICE IELTS ADVANCED READING
% ============================================================
\subsection{Luyện Reading IELTS nâng cao (UC11)}
Bảng~\ref{tab:uc11_spec} đặc tả luồng sự kiện; hình~\ref{fig:activity-uc11} thể hiện sơ đồ activity; hình~\ref{fig:sequence-uc11} thể hiện sơ đồ tuần tự.
\subsubsection{Mô tả use case}
\renewcommand{\arraystretch}{1.3}
\begin{longtable}{|p{0.46\linewidth}|p{0.46\linewidth}|}
    \caption{Đặc tả chức năng Luyện Reading IELTS nâng cao (UC11)} \label{tab:uc11_spec} \\
    \hline
    \multicolumn{2}{|l|}{\textbf{Đặc tả use-case}} \\ \hline
    \multicolumn{2}{|l|}{\textbf{Mã use-case:} UC11} \\ \hline
    \multicolumn{2}{|l|}{\textbf{Actor:} Người học} \\ \hline
    \multicolumn{2}{|p{0.88\linewidth}|}{\textbf{Tiền điều kiện (Pre-condition):} Người học đã đăng nhập thành công và có gói đăng ký Premium hoặc Pro đang hoạt động.} \\ \hline
    \multicolumn{2}{|p{0.88\linewidth}|}{\textbf{Hậu điều kiện (Post-condition):} Phiên luyện tập Reading được lưu vào cơ sở dữ liệu với điểm số, câu trả lời và dữ liệu phân tích theo loại câu hỏi; người học nhận 20~XP; nếu tỷ lệ đúng $\geq 80\%$ nhận thêm 10~XP; chuỗi học liên tiếp được cập nhật.} \\ \hline
    \multicolumn{2}{|l|}{\textbf{Luồng sự kiện chính (Basic flow)}} \\ \hline
    \textbf{Actor} & \textbf{Hệ thống} \\ \hline
    \endfirsthead
    \multicolumn{2}{r}{\textit{(tiếp trang trước)}} \\
    \hline
    \textbf{Actor} & \textbf{Hệ thống} \\ \hline
    \endhead
    \hline
    \multicolumn{2}{r}{\textit{(tiếp trang sau)}} \\
    \endfoot
    \hline
    \endlastfoot
    1. Nhấn vào mục ``IELTS Advanced'' trên thanh điều hướng bên trái (sidebar) của phân hệ IELTS. & \\ \hline
    & 2. Kiểm tra gói đăng ký của người học; hiển thị giao diện IELTS Advanced với 4 tab kỹ năng: Listening, Reading, Writing, Speaking; mặc định đứng ở tab \textbf{Listening}. \\ \hline
    3. Nhấn tab ``Reading''. & \\ \hline
    & 4. Tải danh sách bài luyện tập Reading; hiển thị kèm bộ chọn Part (Part~1, Part~2, Part~3); mặc định chọn Part~1; mỗi thẻ bài kèm tiêu đề, loại câu hỏi và điểm cao nhất đã đạt (nếu đã từng làm). \\ \hline
    5. Nhấn vào một Part (ví dụ: Part~2) trong bộ chọn. & \\ \hline
    & 6. Lọc và hiển thị danh sách các bài luyện tập thuộc Part đã chọn. \\ \hline
    7. Nhấn nút ``Practice Now'' trên một bài luyện tập. & \\ \hline
    & 8. Tải nội dung bài luyện tập; hiển thị giao diện chia hai cột: bên trái là đoạn văn với đánh dấu vị trí câu hỏi, bên phải là các nhóm câu hỏi được render theo dạng bài. \\ \hline
    9. Đọc đoạn văn và điền câu trả lời vào từng câu hỏi. & \\ \hline
    10. Nhấn nút ``Submit Answers'' sau khi hoàn thành. & \\ \hline
    & 11. Chấm điểm từng câu (so sánh không phân biệt hoa thường, hỗ trợ nhiều đáp án chấp nhận được cho Short Answer); tổng hợp kết quả theo từng loại câu hỏi; lưu phiên luyện tập vào cơ sở dữ liệu; kích hoạt sự kiện gamification (20~XP, thêm 10~XP nếu tỷ lệ đúng $\geq 80\%$); cập nhật chuỗi học liên tiếp; điều hướng đến trang Score Report. \\ \hline
    & 12. Hiển thị trang ``Score Report'' gồm: bảng điểm tổng và bảng chi tiết theo loại câu hỏi; phía dưới chia hai cột: bên trái là đoạn văn, bên phải là phần xem lại câu trả lời có hiển thị đáp án; mỗi câu hỏi kèm nút \textbf{Locate}, \textbf{Explain} và \textbf{Note}. \\ \hline
    \multicolumn{2}{|l|}{\textit{Tương tác trên trang Score Report (tùy chọn)}} \\ \hline
    13a. Nhấn nút ``Locate'' trên một câu hỏi. & \\ \hline
    & 14a. Cuộn đoạn văn đến đúng vị trí đánh dấu tương ứng; làm nổi bật vị trí trong khoảng 3~giây rồi tắt. \\ \hline
    13b. Nhấn nút ``Explain'' trên một câu hỏi. & \\ \hline
    & 14b. Mở rộng khung giải thích phía dưới câu hỏi; hiển thị từ khóa, lý giải đáp án và phân tích đáp án nhiễu. \\ \hline
    13c. Nhấn nút ``Note'' trên một câu hỏi để ghi chú cá nhân. & \\ \hline
    & 14c. Hiển thị ô nhập văn bản ghi chú cho câu hỏi tương ứng; người học nhập và lưu nội dung ghi chú. \\ \hline
    \multicolumn{2}{|l|}{\textit{Nhánh B --- Xem lại lịch sử làm bài}} \\ \hline
    15. Nhấn tab ``My Answers'' trên thanh điều hướng của trang bài luyện tập. & \\ \hline
    & 16. Tải lịch sử các phiên làm bài trước đó; hiển thị danh sách gồm ngày giờ, tiêu đề bài và điểm số. \\ \hline
    17. Nhấn ``Details'' trên một phiên làm bài. & \\ \hline
    & 18. Tải dữ liệu phiên đã chọn; hiển thị lại giao diện Score Report với toàn bộ câu trả lời, đáp án và điểm của phiên đó. Kết thúc use-case. \\ \hline
    \multicolumn{2}{|l|}{\textbf{Luồng sự kiện thay thế (Alternative flow)}} \\ \hline
    & 2.1. Người học không có gói Premium hoặc Pro: hệ thống từ chối truy cập và hiển thị màn hình yêu cầu nâng cấp gói đăng ký; use-case kết thúc. \\ \hline
    11.1. Nhấn ``Submit Answers'' khi vẫn còn câu bỏ trống. & \\ \hline
    & 11.2. Nộp bài và chấm điểm; các câu bỏ trống tính là sai. \\ \hline
\end{longtable}

\subsubsection{Sơ đồ activity}
\begin{figure}[H]
    \centering
    \includegraphics[width=\linewidth]{activity-diagram/activity-uc11.png}
    \caption{Sơ đồ activity UC11 -- Luyện Reading IELTS nâng cao}
    \label{fig:activity-uc11}
\end{figure}

\subsubsection{Sơ đồ tuần tự}
\begin{figure}[H]
    \centering
    \includegraphics[width=\linewidth]{sequence-diagram/sequence-uc11.png}
    \caption{Sơ đồ tuần tự UC11 -- Luyện Reading IELTS nâng cao}
    \label{fig:sequence-uc11}
\end{figure}

\subsubsection{Mô tả kỹ thuật chi tiết}

Phân hệ IELTS Advanced Reading được thiết kế theo mô hình giao diện hai cột bất đối xứng, trong đó đoạn văn và bộ câu hỏi được hiển thị đồng thời trên cùng một màn hình. Thiết kế này phản ánh đúng cách người học thực sự làm bài Reading trong kỳ thi IELTS, khi cần đọc đoạn văn và trả lời câu hỏi song song mà không phải cuộn lên xuống giữa hai vùng nội dung.

Kiểm soát truy cập được thực hiện ở tầng máy chủ trước khi trả về bất kỳ nội dung nào: hệ thống xác minh gói đăng ký hiện tại của người học; nếu không đủ điều kiện (không có gói Premium hoặc Pro đang hoạt động), máy chủ từ chối yêu cầu với mã lỗi phù hợp và giao diện hiển thị màn hình yêu cầu nâng cấp. Cơ chế này đảm bảo người dùng không thể tiếp cận nội dung trả phí ngay cả khi can thiệp trực tiếp vào giao diện.

Cấu trúc dữ liệu của bài luyện Reading hỗ trợ đầy đủ các dạng câu hỏi chuẩn IELTS Academic: True/False/Not Given, Matching Headings, Matching Information, Short Answer và Multiple Choice. Mỗi câu hỏi dạng điền từ có thể định nghĩa danh sách đáp án chấp nhận được để xử lý các biến thể chính tả hợp lệ. Phần đoạn văn lưu trữ thông tin vị trí tương ứng của từng câu hỏi trong văn bản, cho phép tính năng ``Locate'' cuộn chính xác đến đoạn chứa đáp án và làm nổi bật tạm thời.

Logic chấm điểm được thực hiện phía máy chủ sau khi người học nộp bài: hệ thống so sánh từng câu trả lời với đáp án chuẩn (không phân biệt hoa thường cho dạng điền từ), tổng hợp kết quả theo từng loại câu hỏi để tạo báo cáo phân tích chi tiết. Kết quả này cùng toàn bộ câu trả lời được lưu vào phiên luyện tập trong cơ sở dữ liệu, cho phép người học xem lại bất kỳ phiên nào trong lịch sử. Sau khi lưu thành công, hệ thống kích hoạt cơ chế gamification và cập nhật chuỗi học liên tiếp trước khi chuyển hướng đến trang báo cáo điểm.

% ============================================================
% UC08 – DO SPEAKING MOCK TEST
% ============================================================
\subsection{Thi thử Speaking (Mock Test) (UC17)}
Bảng~\ref{tab:uc17_spec} đặc tả luồng sự kiện; hình~\ref{fig:activity-uc17} thể hiện sơ đồ activity; hình~\ref{fig:sequence-uc17} thể hiện sơ đồ tuần tự.
\subsubsection{Mô tả use case}
\renewcommand{\arraystretch}{1.3}
\begin{longtable}{|p{0.46\linewidth}|p{0.46\linewidth}|}
    \caption{Đặc tả chức năng Thi thử Speaking (Mock Test) (UC17)} \label{tab:uc17_spec} \\
    \hline
    \multicolumn{2}{|l|}{\textbf{Đặc tả use-case}} \\ \hline
    \multicolumn{2}{|l|}{\textbf{Mã use-case:} UC17} \\ \hline
    \multicolumn{2}{|l|}{\textbf{Actor:} Người học} \\ \hline
    \multicolumn{2}{|p{0.88\linewidth}|}{\textbf{Tiền điều kiện (Pre-condition):} Người học đã đăng nhập thành công.} \\ \hline
    \multicolumn{2}{|p{0.88\linewidth}|}{\textbf{Hậu điều kiện (Post-condition):} Phiên thi Speaking được cập nhật sang trạng thái hoàn thành; kết quả chứa band điểm Speaking ước tính (thang 0.0--9.0) và nhận xét từ AI theo bốn tiêu chí: \textit{Fluency \& Coherence}, \textit{Lexical Resource}, \textit{Grammatical Range \& Accuracy} và \textit{Pronunciation}; chuỗi học liên tiếp được cập nhật.} \\ \hline
    \multicolumn{2}{|l|}{\textbf{Luồng sự kiện chính (Basic flow)}} \\ \hline
    \textbf{Actor} & \textbf{Hệ thống} \\ \hline
    \endfirsthead
    \multicolumn{2}{r}{\textit{(tiếp trang trước)}} \\
    \hline
    \textbf{Actor} & \textbf{Hệ thống} \\ \hline
    \endhead
    \hline
    \multicolumn{2}{r}{\textit{(tiếp trang sau)}} \\
    \endfoot
    \hline
    \endlastfoot
    1. Nhấn vào mục ``IELTS Intensive'' trên thanh điều hướng bên trái của phân hệ IELTS, chọn sub-menu ``Mock Tests''. & \\ \hline
    & 2. Tải danh mục đề thi; hiển thị giao diện Mock Test với bốn tab kỹ năng (Listening, Reading, Writing, Speaking); mặc định đứng ở tab \textbf{Listening}. \\ \hline
    3. Nhấn vào tab ``Speaking''. & \\ \hline
    & 4. Tải danh sách bộ đề Speaking; hiển thị được nhóm theo đầu sách Cambridge IELTS, mỗi thẻ bài kèm số lượng người tham gia, thời gian và band điểm cao nhất đã đạt (nếu đã từng làm). \\ \hline
    5. Nhấn vào một bộ đề Speaking trong danh sách. & \\ \hline
    & 6. Tải thông tin bài thi; hiển thị trang giới thiệu gồm hai phần: \textbf{Exam Details} (thời gian, kỹ năng, cấu trúc 3 Parts và tổng số câu hỏi) và \textbf{Your Examiner} (thông tin giám khảo AI kèm bảng thông tin từng Part theo chuẩn IELTS). \\ \hline
    7. Nhấn nút ``Start Test''. & \\ \hline
    & 8. Hiển thị trang \textbf{Device Test} gồm 3 bước tuần tự --- mỗi bước chỉ được mở khóa khi bước trước hoàn thành: Bước~1 --- Headphone check, Bước~2 --- Microphone check, Bước~3 --- Waiting room. \\ \hline
    9. (Bước 1 --- Headphone check) Nhấn nút phát trên thanh âm thanh để nghe đoạn audio kiểm tra tai nghe. & \\ \hline
    & 10. Phát đoạn audio kiểm tra; thanh tiến trình cập nhật theo thời gian thực; khi audio kết thúc, đánh dấu bước 1 hoàn thành và mở khóa Bước~2. \\ \hline
    11. (Bước 2 --- Microphone check) Nhấn ``REC'', đọc to câu hướng dẫn; nhấn ``STOP'' khi xong; nhấn ``PLAY'' để nghe lại bản ghi âm. & \\ \hline
    & 12. Yêu cầu quyền truy cập microphone; khởi tạo bộ ghi âm và bắt đầu ghi (tự động dừng sau tối đa 10~giây); lưu dữ liệu âm thanh; phát lại khi người học nhấn ``PLAY''; khi phát lại kết thúc, đánh dấu bước 2 hoàn thành và mở khóa Bước~3. \\ \hline
    13. (Bước 3 --- Waiting room) Nhấn nút ``I'M READY''. & \\ \hline
    & 14. Tạo phiên thi mới trong cơ sở dữ liệu với trạng thái đang diễn ra; tải dữ liệu đề thi và phiên; khởi tạo đồng hồ đếm ngược theo thời gian thi; tự động kích hoạt luồng phát câu hỏi đầu tiên của Part~1. \\ \hline
    \multicolumn{2}{|l|}{\textit{--- Speaking Part 1: Introduction \& Interview (lặp lại cho từng câu hỏi) ---}} \\ \hline
    & 15. Hiển thị màn hình Speaking Part~1 với vùng video trung tâm; hiển thị thông báo ``Listen to the question'' trong 2~giây, sau đó tự động phát video câu hỏi của giám khảo. \\ \hline
    & 16. Khi video kết thúc: hiển thị thông báo ``Time to think'' trong 2~giây; kích hoạt bộ đếm thời gian suy nghĩ 2~giây; khi hết bộ đếm, tự động bắt đầu ghi âm câu trả lời. \\ \hline
    17. Nói câu trả lời vào microphone; tùy chọn nhấn nút stop để kết thúc ghi âm sớm. & \\ \hline
    & 18. Ghi âm câu trả lời (tối đa 60~giây, tự động dừng khi hết thời gian); lưu bản ghi âm theo định danh câu hỏi; hiển thị biểu tượng hoàn thành. \\ \hline
    19. Nhấn ``Next''. & \\ \hline
    & 20. Nếu còn câu hỏi tiếp theo trong Part~1: chuyển sang câu kế tiếp, quay lại bước~15. Nếu đã trả lời hết Part~1: chuyển sang giao diện Speaking Part~2. \\ \hline
    \multicolumn{2}{|l|}{\textit{--- Speaking Part 2: Individual Long Turn ---}} \\ \hline
    & 21. Hiển thị giao diện hai cột: cột trái gồm khung video giám khảo và bảng điều khiển ghi âm; cột phải gồm \textbf{Cue Card} (chủ đề và gợi ý nói) và ô ghi chú (\textbf{Note}) cho phép nhập văn bản tự do. Tự động kích hoạt luồng phát câu hỏi. \\ \hline
    & 22. Hiển thị thông báo ``Listen to the question'' (2~giây); phát video giám khảo. Khi video kết thúc: hiển thị thông báo ``Time to think'' (2~giây); kích hoạt bộ đếm thời gian suy nghĩ 60~giây để người học đọc Cue Card và ghi chú. \\ \hline
    23. Đọc Cue Card, ghi chú vào ô Note trong suốt 60~giây thời gian suy nghĩ. & \\ \hline
    & 24. Sau khi hết 60~giây suy nghĩ: nếu câu hỏi Part~2 có video hướng dẫn bổ sung, phát video đó trước khi bắt đầu ghi âm; nếu không, kích hoạt ghi âm trực tiếp. \\ \hline
    25. Nói câu trả lời theo hướng dẫn trong Cue Card; tùy chọn nhấn stop để kết thúc sớm. & \\ \hline
    & 26. Ghi âm câu trả lời (tối đa 120~giây, tự động dừng khi hết thời gian); lưu bản ghi âm; hiển thị biểu tượng hoàn thành. \\ \hline
    27. Nhấn ``Next'' để chuyển sang Part~3. & \\ \hline
    \multicolumn{2}{|l|}{\textit{--- Speaking Part 3: Two-way Discussion (lặp lại cho từng câu hỏi) ---}} \\ \hline
    & 28. Hiển thị màn hình Speaking Part~3; tự động kích hoạt luồng phát câu hỏi theo cùng cơ chế Part~1: caption 2~giây $\to$ phát video $\to$ suy nghĩ 2~giây $\to$ ghi âm tối đa 60~giây. \\ \hline
    29. Nói câu trả lời; tùy chọn nhấn stop; nhấn ``Next'' để chuyển câu. & \\ \hline
    & 30. Nếu còn câu hỏi: chuyển sang câu kế tiếp, quay lại bước~28. Nếu đã trả lời hết tất cả câu hỏi của Part~3: kích hoạt luồng nộp bài. \\ \hline
    \multicolumn{2}{|l|}{\textit{--- Nộp bài và chấm điểm AI ---}} \\ \hline
    31. Nhấn ``Next'' tại câu hỏi cuối của Part~3. & \\ \hline
    & 32. Chuyển sang trạng thái đang nộp bài; hiển thị màn hình chờ AI chấm điểm; tải từng bản ghi âm lên máy chủ và nhận lại đường dẫn công khai; gửi kết quả lên máy chủ kèm danh sách đường dẫn âm thanh và thời gian làm bài. \\ \hline
    & 33. Máy chủ lưu câu trả lời vào phiên thi; đẩy công việc chấm điểm vào hàng đợi xử lý bất đồng bộ. Hệ thống chuyển sang trạng thái đang chấm điểm; định kỳ kiểm tra kết quả chấm điểm; khi AI hoàn thành và trả về kết quả, điều hướng đến trang kết quả. Kết thúc use-case. \\ \hline
    \multicolumn{2}{|l|}{\textbf{Luồng sự kiện thay thế (Alternative flow)}} \\ \hline
    12.1. Từ chối cấp quyền truy cập microphone. & \\ \hline
    & 12.2. Hiển thị cảnh báo yêu cầu cấp quyền microphone; nút ``I'M READY'' không khả dụng; người học không thể tiến hành bài thi. \\ \hline
    & 30.1. Đồng hồ đếm ngược toàn bài về 0 trước khi người học hoàn thành tất cả câu hỏi: bản ghi âm đang thực hiện tự động dừng; hệ thống kích hoạt luồng nộp bài với toàn bộ bản ghi âm đã có; câu hỏi chưa được ghi âm được tính là bỏ trống. \\ \hline
    & 32.1. Lỗi mạng khi tải file âm thanh hoặc khi nộp bài: hệ thống chuyển sang trạng thái lỗi và hiển thị thông báo tương ứng; người học có thể quay lại danh sách đề và thực hiện lại bài thi. \\ \hline
\end{longtable}

\subsubsection{Sơ đồ activity}
\clearpage
\begin{figure}[p]
    \centering
    \includegraphics[width=\linewidth]{activity-diagram/activity-uc17.png}
    \caption{Sơ đồ activity UC17 -- Thi thử Speaking (Mock Test)}
    \label{fig:activity-uc17}
\end{figure}

\subsubsection{Sơ đồ tuần tự}
\begin{figure}[H]
    \centering
    \includegraphics[width=\linewidth]{sequence-diagram/sequence-uc17.png}
    \caption{Sơ đồ tuần tự UC17 -- Thi thử Speaking (Mock Test)}
    \label{fig:sequence-uc17}
\end{figure}

\subsubsection{Mô tả kỹ thuật chi tiết}

\begin{sloppypar}
Phân hệ Speaking Mock Test được triển khai theo mô hình luồng tuần tự nghiêm ngặt, bao gồm ba giai đoạn chính: kiểm tra thiết bị (Device Test), tiến hành thi và nộp bài chấm điểm AI.
\end{sloppypar}

Trong giai đoạn kiểm tra thiết bị, ba bước được liên kết bằng hai cờ trạng thái phía client. Mỗi bước chỉ được kích hoạt khi bước trước hoàn thành --- giao diện ẩn và vô hiệu hóa bước chưa đến lượt để đảm bảo người học không bỏ qua quy trình kiểm tra. Điều kiện hoàn thành bước kiểm tra microphone là người học phải nghe lại bản ghi âm đến cuối, không phải chỉ sau khi ghi xong --- cơ chế này bắt buộc người học xác nhận microphone thực sự hoạt động trước khi bắt đầu thi.

\begin{sloppypar}
Luồng thi Speaking được điều khiển bởi máy trạng thái tám trạng thái: chờ, hiển thị caption, phát video câu hỏi, hiển thị caption suy nghĩ, đang suy nghĩ, phát video thứ hai (chỉ Part~2), đang ghi âm và đã ghi xong. Chuyển trạng thái được kích hoạt bởi sự kiện kết thúc video và các bộ đếm thời gian. Mỗi phần có thông số thời gian khác nhau: Part~1 và Part~3 áp dụng bộ đếm suy nghĩ 2~giây và thời gian ghi âm tối đa 60~giây; Part~2 áp dụng 60~giây suy nghĩ và 120~giây ghi âm.
\end{sloppypar}

Điểm đặc thù của Part~2 là giao diện hai cột: cột phải hiển thị nội dung chủ đề và gợi ý nói (Cue Card) cùng ô ghi chú để người học chuẩn bị trong 60~giây suy nghĩ. Nếu câu hỏi Part~2 có video giám khảo thứ hai, sau khi hết thời gian suy nghĩ hệ thống phát video đó trước khi bắt đầu ghi âm. Thiết kế này mô phỏng đúng quy trình IELTS Speaking thực tế, trong đó giám khảo đọc lại chủ đề trước khi thí sinh bắt đầu nói.

\begin{sloppypar}
Dữ liệu bản ghi âm được lưu phía client theo từng khóa định danh cặp (phần thi -- câu hỏi), cho phép theo dõi trạng thái ghi âm riêng từng câu. Sau khi người học hoàn thành câu hỏi cuối của Part~3, hệ thống thực thi pipeline bất đồng bộ ba bước không chặn giao diện: tải từng bản ghi âm lên máy chủ dưới dạng tệp đính kèm và nhận lại đường dẫn công khai; nộp kết quả kèm danh sách đường dẫn âm thanh và thời gian làm bài; đẩy công việc chấm điểm vào hàng đợi bất đồng bộ. Phía dịch vụ AI, bộ tiêu thụ hàng đợi nhận công việc và sử dụng Gemini AI để đánh giá bốn tiêu chí Speaking, trả về band điểm kèm nhận xét. Hệ thống định kỳ kiểm tra kết quả chấm điểm và chuyển hướng người học đến trang kết quả khi AI hoàn thành, toàn bộ pipeline chạy nền để người học có thể rời trang mà không mất dữ liệu.
\end{sloppypar}

% ============================================================
% UC09 – PRACTICE WRITING WITH DICTATION
% ============================================================
\subsection{Luyện viết với Dictation (UC24)}
Bảng~\ref{tab:uc24_spec} đặc tả luồng sự kiện; hình~\ref{fig:activity-uc24} thể hiện sơ đồ activity; hình~\ref{fig:sequence-uc24} thể hiện sơ đồ tuần tự.
\subsubsection{Mô tả use case}
\renewcommand{\arraystretch}{1.3}
\begin{longtable}{|p{0.46\linewidth}|p{0.46\linewidth}|}
    \caption{Đặc tả chức năng Luyện viết với Dictation (UC24)} \label{tab:uc24_spec} \\
    \hline
    \multicolumn{2}{|l|}{\textbf{Đặc tả use-case}} \\ \hline
    \multicolumn{2}{|l|}{\textbf{Mã use-case:} UC24} \\ \hline
    \multicolumn{2}{|l|}{\textbf{Actor:} Người học} \\ \hline
    \multicolumn{2}{|p{0.88\linewidth}|}{\textbf{Tiền điều kiện (Pre-condition):} Người học đã đăng nhập thành công.} \\ \hline
    \multicolumn{2}{|p{0.88\linewidth}|}{\textbf{Hậu điều kiện (Post-condition):} Tiến trình hoàn thành bài học được lưu vào cơ sở dữ liệu với danh sách các câu đã hoàn thành; trạng thái bài học được cập nhật trong danh sách thư viện; chuỗi học liên tiếp được ghi nhận.} \\ \hline
    \multicolumn{2}{|l|}{\textbf{Luồng sự kiện chính (Basic flow)}} \\ \hline
    \textbf{Actor} & \textbf{Hệ thống} \\ \hline
    \endfirsthead
    \multicolumn{2}{r}{\textit{(tiếp trang trước)}} \\
    \hline
    \textbf{Actor} & \textbf{Hệ thống} \\ \hline
    \endhead
    \hline
    \multicolumn{2}{r}{\textit{(tiếp trang sau)}} \\
    \endfoot
    \hline
    \endlastfoot
    1. Nhấn vào mục ``Shadowing \& Dictation'' trên thanh điều hướng dưới cùng. & \\ \hline
    & 2. Hiển thị màn hình Shadowing \& Dictation với chế độ mặc định là \textbf{Shadowing}; tải danh sách bài học và dữ liệu tiến trình của người học; hiển thị danh sách bài học kèm phần trăm tiến trình Shadowing. \\ \hline
    2. Nhấn vào nút ``Dictation'' trên thanh chuyển chế độ. & \\ \hline
    & 3. Chuyển sang chế độ Dictation; tính lại phần trăm tiến trình theo dữ liệu Dictation của từng bài; hiển thị danh sách bài học kèm thanh tìm kiếm và bộ lọc trạng thái (\textit{All / In Progress / Not Started / Completed}). \\ \hline
    3. (Tùy chọn) Nhập từ khóa vào thanh tìm kiếm hoặc chọn bộ lọc trạng thái. & \\ \hline
    & 4. Lọc danh sách theo tiêu chí: từ khóa khớp tiêu đề bài học và trạng thái tương ứng với phần trăm tiến trình (0\%: Not Started; 1--99\%: In Progress; 100\%: Completed); cập nhật số lượng bài hiển thị. \\ \hline
    4. Nhấn nút ``START'' trên một bài học bất kỳ. & \\ \hline
    & 5. Tải dữ liệu bài học (từ bộ dữ liệu tĩnh đóng gói sẵn hoặc từ máy chủ nếu là video do người dùng thêm); hiển thị màn hình luyện viết gồm phần phát video ở trên và transcript được đục lỗ ở phía dưới; tự động phát đoạn video/audio của câu đầu tiên rồi dừng lại khi hết câu; mức độ khó mặc định là \textbf{Intermediate} (50\% từ bị ẩn). \\ \hline
    5. (Tùy chọn) Chọn mức độ khó trong bộ chọn \textit{Difficulty}. & \\ \hline
    & 6. Áp dụng tỉ lệ đục lỗ tương ứng --- Beginner: hiển thị 70\% từ; Intermediate: 50\%; Advanced: 30\%; Expert: toàn bộ từ bị ẩn; tái tạo tập hợp từ bị ẩn, ưu tiên giữ lại các từ ngắn hơn để lộ trước. \\ \hline
    6. Nghe câu phát ra và điền nội dung vào các ô từ bị đục lỗ trong transcript. & \\ \hline
    7. (Tùy chọn) Nhấn vào một ô từ bị đục lỗ để xem gợi ý. & \\ \hline
    & 8. Hiển thị gợi ý ngay trong ô từ đó: nhấn lần~1 → hiện chữ cái đầu; lần~2 → hiện thêm chữ cái cuối; lần~3 → hiện toàn bộ từ. \\ \hline
    8. (Tùy chọn) Nhấn nút điều chỉnh tốc độ phát. & \\ \hline
    & 9. Chuyển đổi tốc độ theo vòng lặp: $0.25\times \to 0.5\times \to 0.75\times \to 1.0\times \to 2.0\times$; áp dụng ngay cho lần phát tiếp theo. \\ \hline
    9. (Tùy chọn) Nhấn nút phát để nghe lại câu hiện tại. & \\ \hline
    & 10. Seek video đến đầu câu hiện tại và phát lại; thanh tiến trình cập nhật liên tục; tự động dừng khi hết đoạn câu. \\ \hline
    10. Nhấn nút ``Check Answer''. & \\ \hline
    & 11. So sánh từng ô từ đã điền với đáp án (loại bỏ dấu câu, không phân biệt hoa thường); ô đúng chuyển sang nền xanh lá, ô sai chuyển sang nền đỏ. Nếu toàn bộ ô đúng, hiển thị nút ``Next Sentence''. \\ \hline
    11. Nhấn nút ``Next Sentence''. & \\ \hline
    & 12. Đánh dấu câu hiện tại hoàn thành; cập nhật thanh tiến trình tổng thể; chuyển sang câu tiếp theo; tự động phát đoạn video/audio của câu mới; reset trạng thái nhập liệu và tập từ đục lỗ theo mức độ khó hiện tại. Lặp lại từ bước~5. \\ \hline
    12. Sau khi hoàn thành câu cuối cùng, nhấn nút ``Finish ✓''. & \\ \hline
    & 13. Lưu tiến trình vào cơ sở dữ liệu với danh sách các câu đã hoàn thành; hiển thị hộp thoại chúc mừng với hai nút: \textbf{Practice Again} (làm lại từ đầu) và \textbf{Back} (trở về màn hình trước). Kết thúc use-case. \\ \hline
    \multicolumn{2}{|l|}{\textbf{Luồng sự kiện thay thế (Alternative flow)}} \\ \hline
    & 5.1. Bài học không có nguồn video/audio: hiển thị thông báo ``No media available''; người học vẫn thực hiện được dictation dựa trên transcript đục lỗ. \\ \hline
    11.1. Nhấn ``Check Answer'' khi vẫn còn ô từ điền sai. & \\ \hline
    & 11.2. Hiển thị ô đúng màu xanh, ô sai màu đỏ; chỉ hiển thị nút ``Try Again'' thay vì ``Next Sentence''. Quay lại bước 6. \\ \hline
    & 13.1. Lỗi mạng khi lưu tiến trình: hệ thống hiển thị lỗi; người học có thể thử lại hoặc thoát mà không mất kết quả đã thực hiện. \\ \hline
\end{longtable}

\subsubsection{Sơ đồ activity}
\begin{figure}[H]
    \centering
    \includegraphics[width=\linewidth]{activity-diagram/activity-uc24.png}
    \caption{Sơ đồ activity UC24 -- Luyện viết với Dictation}
    \label{fig:activity-uc24}
\end{figure}

\subsubsection{Sơ đồ tuần tự}
\begin{figure}[H]
    \centering
    \includegraphics[width=\linewidth]{sequence-diagram/sequence-uc24.png}
    \caption{Sơ đồ tuần tự UC24 -- Luyện viết với Dictation}
    \label{fig:sequence-uc24}
\end{figure}

\subsubsection{Mô tả kỹ thuật chi tiết}

Màn hình Shadowing \& Dictation được dùng chung cho cả hai chế độ, phân biệt qua biến trạng thái chế độ khởi tạo mặc định là Shadowing. Khi người học chuyển sang Dictation, hệ thống tính lại phần trăm tiến trình của từng bài dựa trên mảng chỉ số các câu đã hoàn thành chia cho tổng số câu của bài đó.

\begin{sloppypar}
Màn hình luyện viết tải bài học theo thứ tự ưu tiên: trước tiên tra cứu danh sách bài học tĩnh đóng gói trong bundle; nếu không tìm thấy, tải từ máy chủ để lấy video do người dùng thêm vào. Việc ưu tiên dữ liệu tĩnh giúp giảm tải mạng và tăng tốc độ khởi tạo màn hình cho phần lớn bài học.
\end{sloppypar}

Cơ chế phát đồng bộ câu hoạt động theo hai nhánh tùy loại nguồn phương tiện: với video YouTube, hệ thống seek đến mốc thời gian bắt đầu câu và sau một khoảng trễ nhỏ bắt đầu phát (để đảm bảo thao tác seek ổn định trước khi phát); với tệp audio thông thường, hệ thống seek và phát ngay. Bộ đếm thời gian theo dõi vị trí phát dựa trên đồng hồ phía client thay vì truy vấn trình phát bất đồng bộ, giúp dừng chính xác khi đến mốc kết thúc câu và mô phỏng đúng hành vi ``phát một câu rồi dừng''.

\begin{sloppypar}
Hệ thống đục lỗ transcript hoạt động theo bốn cấp độ khó ánh xạ sang tỉ lệ từ hiển thị: Beginner (70\%), Intermediate (50\%), Advanced (30\%) và Expert (0\%). Khi cấp độ khó hoặc câu hiện tại thay đổi, hệ thống tái tạo tập hợp từ được lộ: trước tiên loại dấu câu ra khỏi danh sách từ cần ẩn, sau đó sắp xếp các từ còn lại theo chiều dài tăng dần rồi giữ lại đúng số lượng tương ứng với tỉ lệ hiển thị. Chiến lược này ưu tiên lộ các từ ngắn (giới từ, mạo từ) để giảm tải nhận thức cho người học ở cấp độ thấp.
\end{sloppypar}

Kiểm tra đáp án được kích hoạt mỗi khi người học chỉnh sửa nội dung nhập: hệ thống so sánh từng từ đã điền với đáp án chuẩn qua chuẩn hóa văn bản (chuyển về chữ thường, loại dấu câu và khoảng trắng thừa). Khi toàn bộ từ khớp chính xác, hệ thống tự động đánh dấu câu hoàn thành và hiển thị xác nhận. Khi người học hoàn thành bài, hệ thống loại bỏ trùng lặp trong danh sách chỉ số câu đã hoàn thành rồi gửi lên máy chủ. Phía máy chủ, bản ghi tiến trình được tạo mới hoặc cập nhật theo cặp định danh người học -- bài học để đảm bảo mỗi người học chỉ có một bản ghi duy nhất cho mỗi bài.

% ============================================================
% UC10 – ROADMAP
% ============================================================
\subsection{Xây dựng lộ trình học tập (UC19)}
Bảng~\ref{tab:uc19_spec} đặc tả luồng sự kiện; hình~\ref{fig:activity-uc19} thể hiện sơ đồ activity; hình~\ref{fig:sequence-uc19} thể hiện sơ đồ tuần tự.
\subsubsection{Mô tả use case}
\renewcommand{\arraystretch}{1.3}
\begin{longtable}{|p{0.46\linewidth}|p{0.46\linewidth}|}
    \caption{Đặc tả chức năng Xây dựng lộ trình học tập (UC19)} \label{tab:uc19_spec} \\
    \hline
    \multicolumn{2}{|l|}{\textbf{Đặc tả use-case}} \\ \hline
    \multicolumn{2}{|l|}{\textbf{Mã use-case:} UC19} \\ \hline
    \multicolumn{2}{|l|}{\textbf{Actor:} Người học} \\ \hline
    \multicolumn{2}{|p{0.88\linewidth}|}{\textbf{Tiền điều kiện (Pre-condition):} Người học đã đăng nhập thành công và chưa hoàn thành quá trình khởi tạo hồ sơ IELTS.} \\ \hline
    \multicolumn{2}{|p{0.88\linewidth}|}{\textbf{Hậu điều kiện (Post-condition):} Hồ sơ IELTS được tạo mới hoặc cập nhật với mục tiêu band điểm, cam kết học hàng ngày và ngày thi dự kiến; nếu người học làm bài kiểm tra chẩn đoán và đạt điểm $\geq$ 90\% ở kỹ năng nào, toàn bộ bài học kỹ năng đó được đánh dấu hoàn thành trong cơ sở dữ liệu; người học được điều hướng đến lộ trình học tập cá nhân hóa.} \\ \hline
    \multicolumn{2}{|l|}{\textbf{Luồng sự kiện chính (Basic flow)}} \\ \hline
    \textbf{Actor} & \textbf{Hệ thống} \\ \hline
    \endfirsthead
    \multicolumn{2}{r}{\textit{(tiếp trang trước)}} \\
    \hline
    \textbf{Actor} & \textbf{Hệ thống} \\ \hline
    \endhead
    \hline
    \multicolumn{2}{r}{\textit{(tiếp trang sau)}} \\
    \endfoot
    \hline
    \endlastfoot
    1. Nhấn vào mục ``IELTS'' trên thanh điều hướng chính. & \\ \hline
    & 2. Tải trang lộ trình IELTS; phát hiện người học chưa hoàn thành quá trình khởi tạo hồ sơ; tự động điều hướng đến trang onboarding; hiển thị giao diện bước 1 với tiêu đề ``Let's personalize your IELTS roadmap.'' gồm ba mục nhập: \textit{Target IELTS Band Score} (tùy chọn từ 5.0 đến 8.5), \textit{Daily Study Commitment} (15/30/60/120 phút mỗi ngày) và \textit{Target Exam Date} (không bắt buộc). \\ \hline
    2. Chọn band điểm mục tiêu, chọn mức cam kết học hàng ngày; (tùy chọn) chọn ngày thi dự kiến; nhấn ``Continue''. & \\ \hline
    & 3. Lưu tạm thời ba thông số vào trạng thái cục bộ; chuyển sang bước 2; hiển thị giao diện với câu hỏi ``Already know some English?'' và hai lựa chọn: \textbf{Take Short Placement Test} (nút chính) và \textbf{Skip and start from absolute basics} (liên kết phụ). \\ \hline
    3. Nhấn ``Take Short Placement Test''. & \\ \hline
    & 4. Chuyển sang bước 3; tải cặp bài tập chẩn đoán từ máy chủ; hiển thị giao diện bài kiểm tra với thanh tiến trình bốn giai đoạn: \textit{Listening → Reading → Writing → Finish}; chuyển sang giai đoạn Listening. \\ \hline
    & 5. Hiển thị giai đoạn Listening: phát đoạn audio kèm trình phát âm thanh có điều khiển; trình bày 4 câu hỏi dạng trắc nghiệm bên dưới. \\ \hline
    4. Nghe audio, chọn đáp án cho 4 câu hỏi; nhấn ``Submit \& Next''. & \\ \hline
    & 6. Tính điểm Listening theo tỉ lệ phần trăm; cập nhật thanh tiến trình; chuyển sang giai đoạn Reading; hiển thị giao diện hai cột --- bên trái là đoạn văn, bên phải là 2 câu hỏi đọc hiểu. \\ \hline
    5. Đọc passage bên trái, chọn đáp án cho 2 câu hỏi bên phải; nhấn ``Submit \& Next''. & \\ \hline
    & 7. Tính điểm Reading theo cùng công thức; cập nhật thanh tiến trình; chuyển sang giai đoạn Writing; hiển thị giao diện kiểm tra từ vựng dạng điền từ với đoạn văn chứa các ô chỗ trống và danh sách từ chọn. \\ \hline
    6. Chọn từ thích hợp điền vào từng ô chỗ trống trong đoạn văn; nhấn ``Submit \& Next''. & \\ \hline
    & 8. Tính điểm Writing; tính điểm tổng thể trung bình ba kỹ năng; chuyển sang giai đoạn Finish; hiển thị trang ``Diagnostic Complete!'' gồm kết quả từng kỹ năng theo thang chất lượng (Excellent / Good / Needs Work / Beginner) và thông báo lộ trình đã được cá nhân hóa. \\ \hline
    7. Nhấn ``Continue to My Roadmap''. & \\ \hline
    & 9. Gửi thông tin onboarding lên máy chủ bao gồm mục tiêu band điểm, cam kết học, ngày thi và kết quả chẩn đoán; tạo hoặc cập nhật hồ sơ IELTS; với mỗi kỹ năng đạt điểm $\geq$ 90\%, đánh dấu toàn bộ bài học kỹ năng đó là đã hoàn thành; điều hướng người học đến trang lộ trình. Kết thúc use-case. \\ \hline
    \multicolumn{2}{|l|}{\textbf{Luồng sự kiện thay thế (Alternative flow)}} \\ \hline
    3.1. Nhấn ``Skip and start from absolute basics'' (tại bước 2). & \\ \hline
    & 3.2. Tạo hồ sơ IELTS ngay lập tức với điểm chẩn đoán bằng 0, không đánh dấu bài học nào hoàn thành; toàn bộ lộ trình bắt đầu từ bài đầu tiên; điều hướng đến trang lộ trình. \\ \hline
    & 9.1. Lỗi mạng khi lưu hồ sơ: hệ thống hiển thị thông báo lỗi và cho phép người học thử lại mà không mất dữ liệu đã nhập. \\ \hline
\end{longtable}

\subsubsection{Sơ đồ activity}
\begin{figure}[H]
    \centering
    \includegraphics[width=\linewidth]{activity-diagram/activity-uc19.png}
    \caption{Sơ đồ activity UC19 -- Xây dựng lộ trình học tập}
    \label{fig:activity-uc19}
\end{figure}

\subsubsection{Sơ đồ tuần tự}
\begin{figure}[H]
    \centering
    \includegraphics[width=\linewidth]{sequence-diagram/sequence-uc19.png}
    \caption{Sơ đồ tuần tự UC19 -- Xây dựng lộ trình học tập}
    \label{fig:sequence-uc19}
\end{figure}

\subsubsection{Mô tả kỹ thuật chi tiết}

Trang onboarding quản lý luồng ba bước qua biến trạng thái số nguyên. Bước 1 thu thập mục tiêu học tập: người học chọn band điểm trong dải 5.0--8.5, mức cam kết học hàng ngày theo bốn tùy chọn thời gian cố định và ngày thi dự kiến không bắt buộc. Nhấn ``Continue'' chuyển sang bước tiếp theo mà không gọi API, các thông số được giữ trong trạng thái cục bộ.

Luồng phát hiện trạng thái onboarding diễn ra trước khi trang lộ trình IELTS hiển thị: hệ thống kiểm tra hồ sơ IELTS của người dùng; nếu chưa tồn tại hoặc chưa hoàn thành onboarding, máy chủ trả về tín hiệu yêu cầu onboarding và phía client tự động điều hướng đến trang khởi tạo hồ sơ.

\begin{sloppypar}
Bài kiểm tra chẩn đoán được quản lý qua biến trạng thái giai đoạn với bốn bước tuần tự: Listening, Reading, Writing và Finish. Khi khởi tạo, hệ thống tải bài tập Listening và Reading từ máy chủ; giai đoạn Writing sử dụng dữ liệu bài tập tĩnh đóng gói sẵn trong bundle, không cần gọi API riêng. Sau mỗi giai đoạn, điểm số được tính theo tỉ lệ phần trăm và lưu vào trạng thái cục bộ; điểm tổng thể được tính trung bình ba kỹ năng khi hoàn thành giai đoạn Writing.
\end{sloppypar}

Giai đoạn Listening tái sử dụng các thành phần giao diện từ module luyện tập IELTS cơ bản. Giai đoạn Reading bố cục hai cột: đoạn văn bên trái và câu hỏi đọc hiểu bên phải. Giai đoạn Writing dùng dạng bài điền từ: đoạn văn xen kẽ giữa nội dung in sẵn và ô lựa chọn từ, kết xuất liền mạch theo cấu trúc dữ liệu định nghĩa sẵn.

\begin{sloppypar}
Phía máy chủ, quá trình xử lý onboarding sử dụng cơ chế tạo mới hoặc cập nhật để đảm bảo mỗi người dùng chỉ có một hồ sơ IELTS duy nhất, tránh xung đột khi gọi lại nhiều lần. Nếu người học làm bài kiểm tra chẩn đoán, hệ thống duyệt qua ba kỹ năng: với kỹ năng đạt điểm $\geq$ 90\%, toàn bộ bài học và bài tập của kỹ năng đó được đánh dấu hoàn thành trong cơ sở dữ liệu. Cơ chế này đảm bảo người học có trình độ sẵn có không phải lặp lại bài học cơ bản, lộ trình sẽ bắt đầu từ bước phù hợp với năng lực thực tế.
\end{sloppypar}

Sau khi API trả về thành công, giao diện tải lại toàn bộ trang thay vì điều hướng thông thường, vì layout IELTS cần đọc trạng thái hồ sơ mới nhất từ máy chủ để hiển thị đúng nội dung lộ trình sau khi tạo hồ sơ xong.

% ============================================================
% UC11 – POST (Community)
% ============================================================
\subsection{Đăng bài trong Community (UC35)}
Bảng~\ref{tab:uc35_spec} đặc tả luồng sự kiện; hình~\ref{fig:activity-uc35} thể hiện sơ đồ activity; hình~\ref{fig:sequence-uc35} thể hiện sơ đồ tuần tự.
\subsubsection{Mô tả use case}
\renewcommand{\arraystretch}{1.3}
\begin{longtable}{|p{0.46\linewidth}|p{0.46\linewidth}|}
    \caption{Đặc tả chức năng Đăng bài trong Community (UC35)} \label{tab:uc35_spec} \\
    \hline
    \multicolumn{2}{|l|}{\textbf{Đặc tả use-case}} \\ \hline
    \multicolumn{2}{|l|}{\textbf{Mã use-case:} UC35} \\ \hline
    \multicolumn{2}{|l|}{\textbf{Actor:} Người học} \\ \hline
    \multicolumn{2}{|p{0.88\linewidth}|}{\textbf{Tiền điều kiện (Pre-condition):} Người học đã đăng nhập thành công.} \\ \hline
    \multicolumn{2}{|p{0.88\linewidth}|}{\textbf{Hậu điều kiện (Post-condition):} Bài đăng mới được lưu vào cơ sở dữ liệu; bài xuất hiện ngay đầu feed mà không cần tải lại trang; người học nhận +5 XP và huy hiệu bài đăng đầu tiên nếu đây là lần đăng đầu tiên.} \\ \hline
    \multicolumn{2}{|l|}{\textbf{Luồng sự kiện chính (Basic flow)}} \\ \hline
    \textbf{Actor} & \textbf{Hệ thống} \\ \hline
    \endfirsthead
    \multicolumn{2}{r}{\textit{(tiếp trang trước)}} \\
    \hline
    \textbf{Actor} & \textbf{Hệ thống} \\ \hline
    \endhead
    \hline
    \multicolumn{2}{r}{\textit{(tiếp trang sau)}} \\
    \endfoot
    \hline
    \endlastfoot
    1. Nhấn vào mục ``Community'' trên thanh điều hướng. & \\ \hline
    & 2. Tải danh sách bài đăng gần nhất; hiển thị giao diện feed mặc định gồm: ba bộ lọc dạng tag tròn (\textit{All Posts}, \textit{Study Tips}, \textit{Achievements}), ô gợi nhắc ``What's happening?'' kèm thanh công cụ icon và nút ``Post'' ở dưới. \\ \hline
    2. Nhấn vào ô gợi nhắc ``What's happening?'' hoặc nút ``Post'' trên thanh công cụ. & \\ \hline
    & 3. Hiển thị hộp thoại \textbf{Create Post} bao gồm: phần chọn loại bài gồm ba tùy chọn (\textit{GENERAL}, \textit{STUDY TIP}, \textit{SCORE ACHIEVEMENT}) với \textit{GENERAL} được chọn mặc định; ô nhập tiêu đề (tùy chọn); ô văn bản nội dung; mục chọn chủ đề với 9 tag; nút thêm ảnh, nút hủy và nút đăng bài. \\ \hline
    3. Chọn loại bài đăng; (tùy chọn) nhập tiêu đề; nhập nội dung bài đăng; (tùy chọn) nhấn các tag chủ đề phù hợp. & \\ \hline
    & 4. Cập nhật giao diện theo lựa chọn: tag loại được tô màu khi được chọn; tag chủ đề chuyển màu khi được kích hoạt; nút ``Post'' chỉ được kích hoạt khi ô nội dung không rỗng. \\ \hline
    4. Nhấn nút ``Add Image''; chọn một hoặc nhiều file ảnh từ thiết bị. & \\ \hline
    & 5. Tải lên từng ảnh được chọn (tối đa $4 - \text{số ảnh hiện có}$ ảnh, xử lý song song) dưới dạng tệp đính kèm; hiển thị trạng thái đang tải trong khi upload; sau khi thành công, hiển thị lưới xem trước ảnh bên dưới ô nội dung, mỗi ảnh có nút xóa ở góc trên phải. \\ \hline
    5. Nhấn nút ``Post''. & \\ \hline
    & 6. Hiển thị trạng thái đang xử lý; gửi bài đăng lên máy chủ với thông tin loại bài, tiêu đề, nội dung, danh sách chủ đề và đường dẫn ảnh; máy chủ lưu bài vào cơ sở dữ liệu và kích hoạt bất đồng bộ sự kiện gamification (+5 XP); nhận lại bài mới kèm thông tin tác giả; xóa trắng form; đóng hộp thoại; chèn bài mới vào đầu feed ngay lập tức mà không cần tải lại danh sách. Kết thúc use-case. \\ \hline
    \multicolumn{2}{|l|}{\textbf{Luồng sự kiện thay thế (Alternative flow)}} \\ \hline
    3.1. Nhấn ``Cancel'' để hủy bài đăng. & \\ \hline
    & 3.2. Đóng hộp thoại, hủy toàn bộ nội dung đã nhập; feed không thay đổi. \\ \hline
    & 5.1. Đã có đủ 4 ảnh: nút ``Add Image'' bị vô hiệu hóa, người học không thể chọn thêm. \\ \hline
    & 5.2. Lỗi mạng khi tải ảnh lên: hệ thống hiển thị cảnh báo lỗi; các ảnh đã tải thành công trước đó vẫn được giữ nguyên. \\ \hline
    & 6.1. Lỗi mạng khi tạo bài đăng: hệ thống hiển thị cảnh báo lỗi; nội dung form và ảnh đã chọn không bị mất, người học có thể thử lại. \\ \hline
\end{longtable}

\subsubsection{Sơ đồ activity}
\begin{figure}[H]
    \centering
    \includegraphics[width=\linewidth]{activity-diagram/activity-uc35.png}
    \caption{Sơ đồ activity UC35 -- Đăng bài trong Community}
    \label{fig:activity-uc35}
\end{figure}

\subsubsection{Sơ đồ tuần tự}
\begin{figure}[H]
    \centering
    \includegraphics[width=\linewidth]{sequence-diagram/sequence-uc35.png}
    \caption{Sơ đồ tuần tự UC35 -- Đăng bài trong Community}
    \label{fig:sequence-uc35}
\end{figure}

\subsubsection{Mô tả kỹ thuật chi tiết}

Trang Community sử dụng tham số điều hướng để phân biệt các chế độ hiển thị feed: All Posts, Study Tips và Achievements. Ô gợi nhắc đăng bài chỉ hiển thị khi người dùng đã xác thực và tab hiện tại là một trong ba chế độ feed chính, không hiển thị ở tab bảng xếp hạng.

\begin{sloppypar}
Hộp thoại tạo bài đăng là thành phần độc lập quản lý toàn bộ trạng thái form cục bộ: loại bài, tiêu đề, nội dung, danh sách chủ đề, danh sách đường dẫn ảnh đã tải lên và các cờ trạng thái đang tải ảnh/đang gửi bài. Hộp thoại chỉ được kết xuất khi mở, khi đóng sẽ giải phóng hoàn toàn khỏi cây DOM để không chiếm tài nguyên.
\end{sloppypar}

Khi người học chọn ảnh, số lượng ảnh được giới hạn theo số lượng ảnh còn có thể thêm (tối đa 4 ảnh tổng cộng). Các ảnh được tải lên song song, mỗi ảnh gửi dưới dạng tệp đính kèm đa phần đến máy chủ. Phía máy chủ, tệp được tải lên dịch vụ lưu trữ Cloudinary và hệ thống nhận lại đường dẫn công khai dạng HTTPS để lưu vào cơ sở dữ liệu.

\begin{sloppypar}
Trước khi gửi bài, hệ thống kiểm tra nội dung bài đăng không được rỗng ở cả hai tầng: tầng giao diện vô hiệu hóa nút ``Post'' khi ô nội dung rỗng, tầng logic kiểm tra lại trước khi thực hiện yêu cầu. Phía máy chủ, bài đăng được lưu vào cơ sở dữ liệu kèm thông tin tác giả trong cùng một thao tác truy vấn. Sau khi lưu thành công, cơ chế gamification được kích hoạt bất đồng bộ để đảm bảo lỗi trong quá trình xử lý điểm thưởng không ảnh hưởng đến phản hồi trả về cho người dùng.
\end{sloppypar}

Sau khi tạo bài thành công, giao diện chèn bài mới vào đầu danh sách feed hiện tại mà không cần gọi lại API danh sách, tránh độ trễ và giữ nguyên vị trí cuộn của người dùng. Nếu tab hiện tại không phải một trong ba chế độ feed chính, trang tự động điều hướng về feed mặc định để bài mới hiển thị đúng vị trí.

